% Source presentations for register row N-c-016.

\subsection{CTL, CTL*, and state-labelling model checking presentations}
\label{sec:p09-presentation-n-c-016}

\subsubsection{Clarke--Emerson's branching-time synthesis presentation}
\label{sec:p09-presentation-n-c-016-clarke-emerson}

Clarke and Emerson's conference paper appeared in \emph{Logics of Programs},
LNCS~131, on original pages 52--71 in 1982; the meeting itself was held in
1981.  The source used here is the authors' institutionally hosted, authorized
2008 Springer reprint, whose first page identifies those original publication
facts.  Its reprint pagination is 196--215, but citations below use the
original 52--71 pagination.  A different author-hosted scan runs to 37
numbered manuscript pages plus separately scanned figures.  That longer
manuscript can explain what the conference paper calls the \emph{full paper},
but it is not silently treated here as part of the published twenty-page body
\cite[pp.~52--53, 70--71]{ClarkeEmerson1981}.

\paragraph{The program abstraction and its semantic structure.}
The source needs to suppress code details which do not affect
synchronization.  It therefore calls the following per-process object a
\emph{synchronization skeleton}: a finite flowgraph whose uniquely labelled
nodes are regions or blocks of code, whose arcs are possible transitions
between regions, and whose arc labels are the conditions enabling those
transitions.  This is a program abstraction, not a temporal formula or a
semantic state structure.  It differs from an ordinary control-flow graph by
deliberately leaving the computation internal to each region uninterpreted,
and the construction neither synthesizes that suppressed computation nor
proves a refinement from low-level code to the skeleton.  It also differs
from the \emph{global system flowgraph} constructed later:
that object is the finite joint-behavior graph obtained from the synthesized
model by retaining propositional state labels and distinguishing duplicate
labels with auxiliary assignments.  The skeleton belongs to one process,
whereas the global graph records the region of every process.  Unique global
labels are needed to recover unambiguous local enabling conditions, and the
distinction is used when the last synthesis stage factors one global model
into one skeleton per process
\cite[pp.~53--54, 63, 67--70]{ClarkeEmerson1981}.

Fix \(k\in\mathbb N\) with \(k\geq1\) and processes
\(P_1,\ldots,P_k\).  The paper gives the branching-time formulas a semantics
in a structure
\[
 M=(S,A_1,\ldots,A_k,L).
\]
Here \(S\) is countable, \(A_i\subseteq S\times S\) is the transition
relation for one atomic step of process \(P_i\), and \(L(s)\) is the set of
atomic propositions true at \(s\).  Writing
\(A=A_1\cup\cdots\cup A_k\), the source requires \(A\) to be total:
\[
 \forall s\in S\;\exists t\in S\;((s,t)\in A).
\]
A path from \(s_0\) is consequently an infinite sequence
\(s_0,s_1,\ldots\) with \((s_i,s_{i+1})\in A\) for every \(i\).
The computation tree rooted at \(s_0\) is the unfolding of these paths, with
an arc marked by process \(i\) exactly when its underlying transition belongs
to \(A_i\).  Thus the computation tree is neither a per-process skeleton nor
the finite global flowgraph eventually extracted from a model
\cite[pp.~53--54]{ClarkeEmerson1981}.

\paragraph{The source's CTL syntax and semantics.}
The paper calls its language Computation Tree Logic, or CTL, but its primitive
next operator retains the identity of the moving process.  With
\(1\leq j\leq k\), its grammar is
\[
 f ::= p
 \mid \neg f
 \mid (f\land f)
 \mid EX_j f
 \mid A[f\,U\,f]
 \mid E[f\,U\,f].
\]
This is the source's statement family of state formulas.  In particular,
\(EX_j\) is not the unindexed \(EX\) of many later CTL presentations.
Atomic and Boolean satisfaction are classical, while
\[
\begin{aligned}
 M,s\models EX_j f
 &\iff \exists t\in S\,
        \bigl((s,t)\in A_j \land M,t\models f\bigr),\\
 M,s_0\models A[f_1\,U\,f_2]
 &\iff
   \forall(s_0,s_1,\ldots)\,
   \exists i\geq0\,
   \bigl(M,s_i\models f_2
     \land \forall r<i\;M,s_r\models f_1\bigr),\\
 M,s_0\models E[f_1\,U\,f_2]
 &\iff
   \exists(s_0,s_1,\ldots)\,
   \exists i\geq0\,
   \bigl(M,s_i\models f_2
     \land \forall r<i\;M,s_r\models f_1\bigr).
\end{aligned}
\]
The quantifiers in the last two clauses range over paths from \(s_0\).
Because \(i=0\) is permitted, until is inclusive at its right endpoint;
because a witness \(i\) is required, it is strong until
\cite[pp.~54--55]{ClarkeEmerson1981}.

Writing \(\mathit{true}\) for the state formula true at every state, the
source then defines, rather than adds as primitives,
\[
\begin{aligned}
 A[f_1\,V\,f_2]&:=\neg E[\neg f_1\,U\,\neg f_2],
 &
 E[f_1\,V\,f_2]&:=\neg A[\neg f_1\,U\,\neg f_2],\\
 AFf&:=A[\mathit{true}\,U\,f],
 &
 EFf&:=E[\mathit{true}\,U\,f],\\
 AGf&:=\neg EF\neg f,
 &
 EGf&:=\neg AF\neg f,\\
 AX_i f&:=\neg EX_i\neg f,
 &
 EXf&:=\bigvee_{i=1}^{k}EX_i f,
 \qquad AXf:=\neg EX\neg f.
\end{aligned}
\]
The \(V\)-forms are dual release-like operators.  Process-indexed
\(AX_i\) can be vacuously true when process \(i\) has no enabled successor,
even though totality guarantees that some process has a successor.  That
distinction is essential in the mutual-exclusion specification below
\cite[pp.~55--56]{ClarkeEmerson1981}.

\paragraph{The finite-state model checker.}
For a finite \(M\) and an input formula \(f_0\), let
\(\operatorname{sub}^{+}(f_0)\) be the subformulas whose main connective is
not negation.  A state is labelled with exactly one of \(f\) and \(\neg f\)
for each such \(f\), according to its truth.  If
\(n=\operatorname{length}(f_0)\), the source makes \(n+1\) bottom-up passes,
using results for shorter formulas in the pass for formulas of length \(i\).
Atomic, conjunction, and \(EX_j\) cases are immediate.  Until is calculated
by outward radiation from the right-operand states
\cite[pp.~56--58]{ClarkeEmerson1981}.

More explicitly, put
\[
\begin{aligned}
 X_0&=\{s:M,s\models f_2\},&
 X_{r+1}
   &=X_r\cup
     \{s:M,s\models f_1
       \text{ and every }A\text{-successor of }s\text{ lies in }X_r\},\\
 Y_0&=\{s:M,s\models f_2\},&
 Y_{r+1}
   &=Y_r\cup
     \{s:M,s\models f_1
       \text{ and some }A\text{-successor of }s\text{ lies in }Y_r\}.
\end{aligned}
\]
Then \(X_{\lvert S\rvert}\) and \(Y_{\lvert S\rvert}\) are respectively
the sets of states satisfying \(A[f_1\,U\,f_2]\) and
\(E[f_1\,U\,f_2]\).  The \(Y_r\) recurrence is reconstructed from the
paper's printed base case and guarded propagation loop; it is not misreported
as a second displayed recurrence.  Each sequence is monotone over the finite
set \(S\) and therefore stabilizes after at most \(\lvert S\rvert\) strict
additions.  Semantically, a universal until that could postpone \(f_2\) past
every loop-free prefix would admit an infinite path on which \(f_2\) never
occurs; existential until needs only a loop-free witness prefix.  This yields
the paper's
\[
 \operatorname{length}(f_0)\,\lvert S\rvert^2
\]
time bound for basic CTL under its representation and accounting convention.
The later linear edge-sensitive bound below is a different implementation
bound, not a correction to the theorem stated here
\cite[pp.~56--59]{ClarkeEmerson1981}.

There is one typographical defect in the source's displayed approximation.
Its preceding expansion identity, its prose condition, and its printed
Algol-like loop all require the guard \(f_1\):
\[
 A[f_1\,U\,f_2]
 \equiv
 f_2\lor\bigl(f_1\land AX\,A[f_1\,U\,f_2]\bigr).
\]
The intervening display for the \((r+1)\)-st approximation omits
\(f_1\land\).  Reading that display literally would admit a predecessor at
which \(f_1\) is false and would contradict both the semantics and the
algorithm on the next page.  The reconstruction above follows the
proof-supported form and records the omission rather than silently copying it
\cite[pp.~56--57]{ClarkeEmerson1981}.

The paper's worked graph has a root labelled \(a,\neg b\) with successors
labelled \(\neg a,b\) and \(a,\neg b\); the left successor loops, while the
right successor and a state labelled \(a,b\) form a two-cycle.  For
\[
 AFb\land EGa
 \equiv AFb\land\neg AF\neg a,
\]
the \(b\)-states seed \(AFb\).  Backward universal growth next reaches the
right-cycle \(a,\neg b\)-state and then the root, so \(AFb\) holds
everywhere.  By contrast, \(AF\neg a\) grows no farther than the looping
\(\neg a,b\)-state: the all-\(a\) cycle blocks it.  Hence the root receives
\(AFb\) and \(\neg AF\neg a\), exactly the four snapshots in the source's
Figures~4.1--4.4 \cite[pp.~57--58]{ClarkeEmerson1981}.

The paper separately extends its notation with \(E^{\infty}Fp\), meaning
that some path visits \(p\)-states infinitely often.  Proposition~4.1 proves
that, in a finite structure, this holds at \(s\) exactly when some reachable
\(p\)-state \(s'\) either has a self-loop or belongs to a strongly connected
component of cardinality greater than one.  Reachability plus one graph
search therefore preserves the same displayed quadratic bound.  This
infinitary addition is not one of the basic grammar clauses above
\cite[pp.~57--59]{ClarkeEmerson1981}.

Nor does the paper claim that arbitrary linear-time path assertions retain
that bound.  Proposition~4.2 prefixes an LTL assertion by a path quantifier
and reduces directed Hamiltonian path to model checking.  Given a directed
graph \(G=(V,E)\) with \(n\geq1\) and
\(V=\{v_1,\ldots,v_n\}\), retain every edge in \(E\), add fresh states
\(u_1,u_2\), and add exactly the new edges
\[
 u_1\to v_i,\qquad v_i\to u_2\quad(1\leq i\leq n),
 \qquad u_2\to u_2.
\]
Label \(v_i\) only by \(p_i\), and give \(u_1,u_2\) none of the \(p_i\).
At \(u_1\), the source tests the formula
\[
 E\left[
   \bigwedge_{i=1}^{n}Fp_i
   \land
   \bigwedge_{i=1}^{n}G\bigl(p_i\to XG\neg p_i\bigr)
 \right].
\]
Any satisfying path from \(u_1\) must visit every labelled vertex and may
visit none twice before remaining at \(u_2\); it therefore encodes a directed
Hamiltonian path in \(G\), and every such Hamiltonian path extends through
\(u_2\) to a satisfying infinite path.  The literal conclusion is
\(\mathsf{NP}\)-hardness
for this extended model-checking problem, not an \(\mathsf{NP}\)-completeness
classification and not the later full CTL* theorem
\cite[pp.~59--60]{ClarkeEmerson1981}.

\paragraph{What the satisfiability section actually supplies.}
Section~5 expressly calls its tableau decision procedure an \emph{outline}.
Its announced input/output contract is: on \(f_0\), return satisfiable or
unsatisfiable and, in the satisfiable case, construct a finite model.  The
three advertised stages are to build a finite directed AND/OR tableau, delete
its inconsistent portions, and unravel what remains into a model.  The
twenty-page paper does not print a numbered soundness-and-completeness theorem
or its full proof; it directs the reader to the \emph{full paper}.  Its
footnote also says that the related \([BM81]\) algorithm is incorrect because
it can reject satisfiable formulas, and it cites corrected procedures instead.
Those limitations belong to the source's proof status
\cite[pp.~60--63]{ClarkeEmerson1981}.

The source uses \(G\) for an OR-node label and \(H\) for an AND-node label.
The publication names two successor constructions on finite formula sets.  It
defines \(\operatorname{Blocks}(G)\) through the local Boolean and fixed-point
closure below, but it illustrates \(\operatorname{Tiles}(H)\) only through the
printed example below rather than giving the full general process-indexed
clauses.  Their announced roles are nonredundant: a block chooses one locally
satisfiable alternative, while every successor obligation represented by the
printed tiles must be satisfiable.  The publication alternates the two named
constructions in its AND/OR-tableau outline:
\[
 G\longrightarrow\operatorname{Blocks}(G)
 \longrightarrow\operatorname{Tiles}(H)
 \longrightarrow\operatorname{Blocks}(G')
 \longrightarrow\cdots.
\]

For the first operation, an \(\alpha\)-formula has two required components
and a \(\beta\)-formula has two alternative components.  Thus
\(f\land g\) is \(\alpha\), \(f\lor g\) is \(\beta\), and the fixed-point
expansions make
\[
 EFp\equiv p\lor EXEFp
 \quad\text{a }\beta\text{-formula},\qquad
 AGp\equiv p\land AXAGp
 \quad\text{an }\alpha\text{-formula}.
\]
A formula with no modality, or with main connective \(EX\) or \(AX\), is
called \emph{elementary} and is not expanded further.  Every other formula is
nonelementary; the stated \(\alpha/\beta\) component rules govern those
nonelementary expansions.  A finite formula
set is downward closed when it contains both components of every
\(\alpha\)-formula it contains and at least one component of every
\(\beta\)-formula it contains.  The members
\(H\in\operatorname{Blocks}(G)\) are precisely the downward-closed results
of repeatedly making those expansions; \(G\) is satisfiable exactly when
some such \(H\) is satisfiable.

For the second operation, suppose the elementary part of \(H\) is
\[
 \{AXh_1,AXh_2,EXg_1,EXg_2,EXg_3\}.
\]
Then the source's example gives
\[
 \operatorname{Tiles}(H)
 =
 \bigl\{
   \{h_1,h_2,g_1\},
   \{h_1,h_2,g_2\},
   \{h_1,h_2,g_3\}
 \bigr\}.
\]
Each existential next formula needs a witness successor; every such witness
must also meet both universal next obligations.  Provided \(H\) has no
propositional inconsistency, \(H\) is satisfiable exactly when all of these
tiles are satisfiable.  Labels of the same node type are merged.  Since there
are at most \(2^{\operatorname{length}(f_0)}\) labels and each may occur once
as each of the two node types, construction terminates
\cite[pp.~60--62]{ClarkeEmerson1981}.

The worked tableau begins with
\(G_0=\{EFp\land EF\neg p\}\).  Boolean/fixed-point expansion produces four
candidate AND-nodes, including one containing both \(p\) and \(\neg p\),
which is deleted.  Consistency deletion also requires each eventuality
\(EFh\) or \(AFh\) to have a reachable node where \(h\) holds.  In the
remaining tableau the root is not deleted.  Choosing the root copy where
\(p\) already fulfills \(EFp\), then its successor copy where
\(\neg p\) fulfills \(EF\neg p\), yields a finite model; the latter state can
loop.  This example illustrates the announced construction, but it does not
turn the abridged outline into a printed full correctness proof
\cite[pp.~61--63]{ClarkeEmerson1981}.

\paragraph{Synthesis and the mutual-exclusion derivation.}
The synthesis method has exactly three stages: write the required concurrent
behavior in the source's CTL; apply the satisfiability procedure to obtain a
finite global model; then factor the individual synchronization skeletons
from that model's global system flowgraph.  The central example has two
processes, each with the atomic region labels
\(\mathsf{NCS}_i\), \(\mathsf{TRY}_i\), and \(\mathsf{CS}_i\) for its
noncritical, trying, and critical regions.  In the following nine requirement
groups, \(i\in\{1,2\}\) and \(j\) is the other process.  Their conjunction is
the formula passed to the tableau:
\begin{enumerate}
\item The designated start state is
\[
 \mathsf{NCS}_1\land\mathsf{NCS}_2.
\]
\item Mutual exclusion is
\[
 AG\neg(\mathsf{CS}_1\land\mathsf{CS}_2).
\]
\item Absence of starvation for each process is
\[
 AG(\mathsf{TRY}_i\to AF\,\mathsf{CS}_i).
\]
\item Exactly one region is occupied:
\[
\begin{gathered}
 AG(\mathsf{NCS}_i\lor\mathsf{TRY}_i\lor\mathsf{CS}_i),\\
 AG(\mathsf{NCS}_i\to
     \neg(\mathsf{TRY}_i\lor\mathsf{CS}_i)),\\
 AG(\mathsf{TRY}_i\to
     \neg(\mathsf{NCS}_i\lor\mathsf{CS}_i)),\\
 AG(\mathsf{CS}_i\to
     \neg(\mathsf{NCS}_i\lor\mathsf{TRY}_i)).
\end{gathered}
\]
\item Entry from the noncritical region into the trying region remains
possible:
\[
 AG(\mathsf{NCS}_i\to EX_i\mathsf{TRY}_i).
\]
\item Every enabled move by process \(i\) from its trying region enters its
critical region:
\[
 AG\bigl(
   \mathsf{TRY}_i\land EX_i\mathit{true}
   \to AX_i\mathsf{CS}_i
 \bigr).
\]
The existential conjunct prevents the universal next clause from satisfying
the requirement merely because process \(i\) has no enabled move.
\item Return from the critical region to the noncritical region remains
possible:
\[
 AG(\mathsf{CS}_i\to EX_i\mathsf{NCS}_i).
\]
\item A step by \(j\) cannot change the region occupied by \(i\):
\[
\begin{gathered}
 AG(\mathsf{NCS}_i\to AX_j\mathsf{NCS}_i),\\
 AG(\mathsf{TRY}_i\to AX_j\mathsf{TRY}_i),\\
 AG(\mathsf{CS}_i\to AX_j\mathsf{CS}_i).
\end{gathered}
\]
\item Some process can always move:
\[
 AG(EX\mathit{true}).
\]
\end{enumerate}
Under the standing totality condition on
\(A=A_1\cup\cdots\cup A_k\), group~9 is semantically redundant: every state
already has an \(A_i\)-successor for some \(i\).  It is retained because it is
one of the source's printed requirement groups.
These are formulas in the source's indexed logic, not informal requirements
retrofit to a later unindexed CTL grammar
\cite[pp.~63--65]{ClarkeEmerson1981}.

The tableau is simplified by propagating the consequences of invariants
without copying every invariant into every label and by suppressing
propositional consequences already forced by a region label.  It deletes,
among other nodes, the candidate containing both
\(\mathsf{CS}_1\) and \(\mathsf{CS}_2\).  Here a \emph{model fragment} is a
rooted finite directed acyclic graph of AND-nodes embeddable in the tableau
and chosen so that every eventuality at its root is fulfilled inside it.  It
is a construction artifact, not a complete model; its root supplies the
checkpoint which permits later cycle closure without abandoning an
eventuality.  The global model is assembled from these fragments, and cycles
are closed only through checked fragment roots.  Its states initially
carry the pair of current process regions.  Two model nodes nevertheless have
the same pair
\((\mathsf{TRY}_1,\mathsf{TRY}_2)\).  The source introduces
\(\mathsf{TURN}\) solely to distinguish those nodes: incoming arcs assign
\(\mathsf{TURN}:=1\) or \(\mathsf{TURN}:=2\), and the corresponding equality
is added to the target label.  This auxiliary program variable is a component
of the global-state label, not a temporal connective.  It is nonredundant
because the factorization rule below requires global nodes to have distinct
labels, and it is used in the enabling conditions of the resulting skeletons
\cite[pp.~65--68]{ClarkeEmerson1981}.

For example, the factored skeleton for \(P_1\) has the following
synchronization-relevant transitions:
\[
\begin{array}{rcll}
 \mathsf{NCS}_1&\longrightarrow&\mathsf{TRY}_1
   &\text{when }\mathsf{NCS}_2\lor\mathsf{CS}_2,\\
 \mathsf{NCS}_1&\longrightarrow&\mathsf{TRY}_1
   &\text{when }\mathsf{TRY}_2,
     \text{ assigning }\mathsf{TURN}:=2,\\
 \mathsf{TRY}_1&\longrightarrow&\mathsf{CS}_1
   &\text{when }\mathsf{NCS}_2
     \lor(\mathsf{TRY}_2\land\mathsf{TURN}=1),\\
 \mathsf{CS}_1&\longrightarrow&\mathsf{NCS}_1
   &\text{when }\mathsf{NCS}_2\lor\mathsf{TRY}_2.
\end{array}
\]
The skeleton for \(P_2\) is obtained by the source's displayed symmetry,
interchanging process indices and assigning \(\mathsf{TURN}:=1\) when
\(P_2\) starts trying while \(P_1\) is already trying.  The two
\(\mathsf{TRY}\)-states with different turn values thereby enable different
critical-section entries; the auxiliary value is not decorative
\cite[pp.~67--69]{ClarkeEmerson1981}.

The general factorization makes this calculation precise.  Retain only the
propositional labels of model nodes.  The reader-facing typed form of the
source's auxiliary-value step is as follows.  If one such label \(\Lambda\)
occurs on distinct nodes \(v_1,\ldots,v_d\), where \(d\geq2\), introduce a
fresh auxiliary program variable \(z_{\Lambda}\) with value domain
\(\{0,1,\ldots,d\}\).  At every global node \(w\), add exactly one state
proposition for this variable:
\[
 z_{\Lambda}=r\quad\text{if }w=v_r,
 \qquad
 z_{\Lambda}=0\quad\text{if }w\notin\{v_1,\ldots,v_d\}.
\]
Initialize every \(z_{\Lambda}\) at every designated initial node to the value
shown by that node's proposition, regardless of its incoming arcs.  On every
global arc, assign each \(z_{\Lambda}\) the value shown at the target node.
Thus every fresh variable is total, \(\Lambda\) remains a propositional label,
and \(z_{\Lambda}\) records exactly which member of its duplicate class is
current, or \(0\) outside that class.  The result is the \emph{global system
flowgraph}: a finite, uniquely labelled joint-behavior graph.  For process
\(P_i\) with \(m\geq1\) regions \(R_1,\ldots,R_m\), begin its
skeleton with those region nodes and no arcs.  This rule applies only when
every global \(i\)-arc preserves every other process's region component:
\[
 \forall j\ne i,\qquad
 \text{the \(P_j\)-region at the target equals its region at the source}.
\]
Requirement group~8 establishes this premise for the displayed
mutual-exclusion construction.  A global arc violating it lies outside this
factorization rule because a local \(P_i\) transition cannot update a
\(P_j\)-region.  For a global \(i\)-arc \(e\)
whose \(P_i\)-component changes from \(R_u\) to \(R_v\), let \(\gamma_e\) be
the conjunction of every conjunct in the global source label except the atom
asserting \(P_i\)'s own current region.  Thus \(\gamma_e\) contains the
\(k-1\) other-process region components and every auxiliary equality true
at the source node.  Each empty part of this conjunction is
\(\mathit{true}\).  Let \(\sigma_e\) be the arc's complete auxiliary update
map, with the absence of an assignment to a variable represented by the
identity update for that variable.  Here \(\gamma_e\) is a state-formula guard
and \(\sigma_e\) is a transition update; they are metalinguistic
abbreviations with different standings.
For every \(u,v\in\{1,\ldots,m\}\) and every update map \(\sigma\) realized
by at least one such global arc, create one local transition
\(R_u\to R_v\), attach update \(\sigma\), and give it the guard
\[
 \bigvee\left\{
   \gamma_e
   \;\middle|\;
   e\text{ is a global }i\text{-arc from }R_u\text{ to }R_v
   \text{ and }\sigma_e=\sigma
 \right\}.
\]
Only arcs with the same complete update map are merged.  Distinct updates,
including identity versus a nonidentity assignment, remain parallel local
transitions, exactly as in the two displayed
\(\mathsf{NCS}_1\to\mathsf{TRY}_1\) branches.  This is the point at which a
finite global model becomes executable per-process enabling conditions
\cite[pp.~68--70]{ClarkeEmerson1981}.

The source's conclusion is deliberately provisional.  It says the synthesis
algorithm may have exponential behavior, argues only that small skeletons
might make this manageable, and calls practical feasibility an open research
problem.  Its empirical claim is limited to an experimental C implementation
of the model checker on a DEC~11/70; that implementation is not a proof of
the synthesis procedure.  Accordingly, this presentation preserves four
different statuses: a printed finite-state labelling algorithm with its source
argument and bound; a proposition-level SCC characterization and
\(\mathsf{NP}\)-hardness reduction with printed proofs; an abridged tableau
decision-procedure outline whose details the publication refers to the full
paper; and a worked synthesis/factorization construction for which this
publication states no separate general correctness theorem
\cite[pp.~59--63, 69--70]{ClarkeEmerson1981}.

\subsubsection{Queille--Sifakis's CESAR presentation}
\label{sec:p09-presentation-n-c-016-queille-sifakis}

The controlling publication is the fifteen-page LNCS chapter of 1982.
Springer's 2008 reprint preserves it as pp.~216--230.  The complete original
body was read in a publicly accessible institutional scan and authenticated
against Springer's original and reprint previews; the scan's distribution
authorization is not established, and the subscription-only reprint pages
218--230 were not independently inspected.  Those rights and version limits
affect access provenance, not the textual identity of the reconstructed
original body \cite{QueilleSifakis1982}.

\paragraph{CESAR's program model and its exact output claim.}
CESAR takes two unlike inputs: a high-level description program for
communicating sequential processes and branching-time specification formulas.
The translator maps the program to an \emph{interpreted Petri net} (IPN):
a Petri net equipped with a vector \(\mathbf X\) of the program variables
retained in the IPN and, for each transition \(t_i\), a guarded command
\[
 c_i(\mathbf X)\longrightarrow \mathbf X:=\alpha_i(\mathbf X).
\]
The translator type-checks and deletes internal variables of non-specified
types; such deleted source-program variables are not components of
\(\mathbf X\).
This source term is necessary because an ordinary Petri net records only
control-token flow, whereas CESAR's transition also tests and changes data.
Its standing is the program model analyzed by the system, not a
formula or a proof.  It differs from an ordinary net exactly by those guards
and assignments, and it is used to calculate the predecessor transformer
below.  A transition fires only when both its Petri-net enabling condition and
\(c_i\) hold, and firing performs \(\alpha_i\).  A missing inscription means
the true guard and identity assignment.

Processes communicate by rendezvous on exchanged variables: an output
\(!V:=e\) in one process synchronizes with an input \(?V\) in another.
The translator constructs one safe state graph per sequential process and
merges rendezvous transitions.  The paper asserts preservation of each
process's safety, but sends the graph-grammar translation elsewhere and does
not prove translation correctness here.  Figure~1 names the final system
output only as \emph{diagnostics}; it supplies no counterexample format,
synthesis algorithm, or low-level refinement theorem
\cite[pp.~337--340]{QueilleSifakis1982}.

The complete example is the alternating-bit protocol.  In this example,
scalar \(X\) is the sender's data value, not the IPN variable vector
\(\mathbf X\); \(Y\) and \(Z\) are respectively the sender's and receiver's
control bits.  Its four processes
are a sender, a receiver, a sender-to-receiver line, and a
receiver-to-sender line.  The sender transmits \((X,Y)\), accepts only an
acknowledgement carrying its current bit \(Y\), and then complements \(Y\).
The receiver accepts only a message carrying its current bit \(Z\), sends the
matching acknowledgement, and then complements \(Z\).  Each line can either
transmit its value or take a true-guarded loss branch.  Figure~2 is the
translated four-component safe IPN, with initially marked places
\[
 s_1,\quad r_1,\quad m_1,\quad a_1
\]
and the merged rendezvous transitions, guards, and assignments.

The outgoing message is \(M=(X,Y)\), \(MM\) is its
receiver-side copy, \(AA\) is the outgoing acknowledgement, and the field
projection \((\cdot).B\) returns a message or acknowledgement bit.  The
identifiers \(\mathsf{SENDER}\) and \(\mathsf{RECEIVER}\) name the first two
processes.  The action labels used below have these exact program-event
standings:
\[
\begin{array}{c|l}
\mathit{send}&\text{initial output of }M:=(X,Y)\\
\mathit{repeat}&\text{retransmission of that same message}\\
\mathit{receive}&\text{receiver input of }MM\\
\mathit{accept}&\text{receiver acceptance when }MM.B=Z\\
\mathit{repeatack}&\text{retransmission of the previous acknowledgement}\\
\mathit{sendack}&\text{acknowledgement output after acceptance}\\
\mathit{receiveack}&\text{sender input of an acknowledgement}
\end{array}
\]
The labels are noninterchangeable because the \(\operatorname{after}\) and
\(\operatorname{enable}\) formulas below distinguish these events.

The prose assumes that the line cannot lose every message and
acknowledgement.  The displayed programs do not encode that assumption:
both line processes may choose the loss branch forever, and the temporal
semantics below has no fairness clause.  The assumption is therefore external
to the printed transition model.  Nothing proved from the program alone may
silently use it \cite[pp.~340--343]{QueilleSifakis1982}.

\paragraph{The literal branching-time language.}
Let \(F\) be the atomic state predicates.  The source language is
\[
\begin{aligned}
 f::={}&p\mid\mathsf{true}\mid\mathsf{false}\mid
 \neg f\mid f\land f\mid f\lor f\mid f\Rightarrow f\\
 &\mid\operatorname{POT}(f)\mid\operatorname{INEV}(f),
 \qquad p\in F,
\end{aligned}
\]
with the defined duals
\[
 \operatorname{ALL}(f):=\neg\operatorname{POT}(\neg f),
 \qquad
 \operatorname{SOME}(f):=\neg\operatorname{INEV}(\neg f).
\]
These are unary state-formula operators; the source has no separate
path-formula grammar, next operator, or until operator.  In later CTL
terminology their behaviors are respectively \(EF,AF,AG,EG\), but the 1982
language is not thereby the full later CTL grammar.

Write the transition structure as \(S=(Q,\to)\).  An execution from \(q\) is
a finite or infinite sequence \(s\) with \(s(0)=q\) and
\(s(i)\to s(i+1)\) whenever the next position is defined.  It is maximal when
it is infinite or finite at a sink.  These start and adjacency clauses are
this presentation's explicit completion of the printed omission recorded
below.  For an execution \(s\), let
\[
 D_s:=\{k\in\mathbb N:s(k)\text{ is an actual position of }s\}.
\]
Thus \(D_s=\mathbb N\) for an infinite execution and
\(D_s=\{0,\ldots,n\}\) for a finite execution
\(s(0),\ldots,s(n)\).  The paper adjoins a fictitious inaccessible state
\(\omega\) after a finite execution ends, and its indexed reachability glyph
satisfies \(q\xrightarrow{k}s(k)\) exactly when \(k\in D_s\).  Because formula
satisfaction is defined on \(Q\), not on \(\omega\), the total typed form of
the temporal clauses quantifies only over \(D_s\).  With \(EX_q\) the maximal
executions from \(q\), they are
\[
\begin{aligned}
 q\models\operatorname{POT}(f)
 &\Longleftrightarrow
 \exists s\in EX_q\;\exists k\in D_s\;(s(k)\models f),\\
 q\models\operatorname{INEV}(f)
 &\Longleftrightarrow
 \forall s\in EX_q\;\exists k\in D_s\;(s(k)\models f),\\
 q\models\operatorname{ALL}(f)
 &\Longleftrightarrow
 \forall s\in EX_q\;\forall k\in D_s\;(s(k)\models f),\\
 q\models\operatorname{SOME}(f)
 &\Longleftrightarrow
 \exists s\in EX_q\;\forall k\in D_s\;(s(k)\models f).
\end{aligned}
\]
All four permit the current state as \(k=0\).  Restriction to \(D_s\) is the
typed form of the source's reachability guard and never asks the fictitious
post-terminal state to satisfy a formula.

The printed definition is formally incomplete at two places.  Its sentence
introducing an execution from \(q\) omits both \(s(0)=q\) and the adjacency
clauses \(s(i)\to s(i+1)\).  It also names
\(\mathsf{false},\lor,\Rightarrow\) but prints interpretation clauses only
for \(\mathsf{true},\neg,\land\) and the temporal operators.  The intended
ordinary clauses are determinate from use, but they are not misreported as
displayed source text \cite[pp.~344--345]{QueilleSifakis1982}.

\paragraph{Program predicates and the complete property families.}
For a translated program, \(\operatorname{Init}\) is the initial-state
predicate.  For a labelled action \(a\),
\(\operatorname{enable}a\) holds exactly where \(a\) can execute, and
\(\operatorname{after}a\) holds immediately after it terminates.  These are
atomic state predicates derived from the program/IPN, not temporal operators.
Their distinction is necessary because an action may be enabled without
being selected and a state after an action need not retain that enabledness.
They are used in every source specification below.

For every \(k\in\mathbb N\) and finite action family
\((a_i)_{i=1}^{k}\), the source's finite-family abbreviations are
\[
\begin{aligned}
 \operatorname{enable}(a_1,\ldots,a_k)
   &:=\bigvee_{i=1}^{k}\operatorname{enable}a_i,\\
 \operatorname{after}(a_1,\ldots,a_k)
   &:=\bigvee_{i=1}^{k}\operatorname{after}a_i.
\end{aligned}
\]
For a process \(\mathcal P\) with finite action set \(A(\mathcal P)\), the
source also writes \(\operatorname{enable}\mathcal P\) for the actions in that
set.  Its reader-facing eliminable value is
\[
 \operatorname{enable}\mathcal P
 :=\bigvee_{a\in A(\mathcal P)}\operatorname{enable}a.
\]
For the reader-facing finite-family notation, adopt the nullary conventions
\[
 \bigvee\varnothing=\mathsf{false},
 \qquad
 \bigwedge\varnothing=\mathsf{true}.
\]
Thus at \(k=0\) both finite-family abbreviations are false, and a process with
\(A(\mathcal P)=\varnothing\) has
\(\operatorname{enable}\mathcal P=\mathsf{false}\); the conjunction saying
that every member of an empty family is live is true.  For a Boolean program
predicate \(P\), the source's invariant schema is
\[
 \operatorname{Init}\Rightarrow\operatorname{ALL}(P).
\]
For the protocol, the three printed instances say that a sent message carries
the sender bit, the line preserves a received message, and the received bit
equals the sender bit:
\[
\begin{aligned}
\operatorname{Init}&\Rightarrow
 \operatorname{ALL}(
   \operatorname{after}(\mathit{send},\mathit{repeat})\Rightarrow M.B=Y),\\
\operatorname{Init}&\Rightarrow
 \operatorname{ALL}(\operatorname{after}\mathit{receive}\Rightarrow MM=M),\\
\operatorname{Init}&\Rightarrow
 \operatorname{ALL}(\operatorname{after}\mathit{receive}\Rightarrow MM.B=Y).
\end{aligned}
\]

The source calls \(a\) live when it is possible to reach an enabling state
from every reachable state:
\[
 \operatorname{Init}\Rightarrow
 \operatorname{ALL}\operatorname{POT}\operatorname{enable}a.
\]
For a set of actions, liveness of every member uses a conjunction of those
targets, whereas absence of deadlock uses their disjunction:
\[
\begin{aligned}
\operatorname{Init}&\Rightarrow
 \operatorname{ALL}\left(
   \bigwedge_{i=1}^{k}\operatorname{POT}\operatorname{enable}a_i\right),\\
\operatorname{Init}&\Rightarrow
 \operatorname{ALL}\operatorname{POT}
   \operatorname{enable}(a_1,\ldots,a_k).
\end{aligned}
\]
The four protocol targets are, in increasing source strength,
\(\operatorname{enable}(\mathsf{SENDER})\),
\(\operatorname{enable}(\mathsf{RECEIVER})\),
\(\operatorname{enable}\mathit{send}\), and
\(\operatorname{enable}\mathit{accept}\).  These formulas assert possible
future enabledness; they neither assert occurrence on every execution nor add
fairness.

The two response schemas are
\[
\begin{aligned}
\operatorname{Init}&\Rightarrow
 \operatorname{ALL}(\operatorname{after}a\Rightarrow
   \operatorname{POT}\operatorname{enable}b),\\
\operatorname{Init}&\Rightarrow
 \operatorname{ALL}(\operatorname{after}a\Rightarrow
   \operatorname{INEV}\operatorname{enable}b).
\end{aligned}
\]
Again the conclusion is that \(b\) becomes executable, not that it executes.
The first three printed protocol instances use:
\[
\begin{array}{c|c|c}
\text{trigger}&\text{temporal demand}&\text{target}\\ \hline
\operatorname{after}(\mathit{send},\mathit{repeat})&\operatorname{POT}
  &\operatorname{enable}\mathit{receive}\\
\operatorname{after}(\mathit{sendack},\mathit{repeatack})&\operatorname{POT}
  &\operatorname{enable}\mathit{receiveack}\\
\operatorname{after}(\mathit{send},\mathit{repeat})&\operatorname{INEV}
  &\operatorname{enable}\mathit{receiveack}
\end{array}
\]
The fourth is the single formula
\[
 \operatorname{Init}\Rightarrow
 \operatorname{ALL}\Bigl(
   (\operatorname{after}\mathit{accept}\Rightarrow
     \operatorname{INEV}(\operatorname{enable}\mathit{sendack}))
   \land
   (\operatorname{after}\mathit{sendack}\Rightarrow AA.B=MM.B)
 \Bigr).
\]
In the second and third displayed source formulas, the identifier is
misprinted as \(\operatorname{receivack}\); the program and Figure~2 define
only \(\operatorname{receiveack}\).  The reader-facing table repairs the
identifier while recording the defect
\cite[pp.~345--347]{QueilleSifakis1982}.

\paragraph{The four fixed-point calculations.}
Identify a predicate with its subset of \(Q\), and write
\(\lvert f\rvert=\{q\in Q:q\models f\}\).  Definition~1 gives the
existential predecessor transformer
\(\operatorname{pre}:\mathcal P(Q)\to\mathcal P(Q)\),
\[
 \operatorname{pre}(P)
 =\{q\in Q:\exists q'\,(q\to q'\land q'\in P)\},
\]
and its Boolean dual is
\[
 \widetilde{\operatorname{pre}}(P)
 =Q\setminus\operatorname{pre}(Q\setminus P).
\]
The first asks for some immediate successor in \(P\); the dual says that no
immediate successor lies outside \(P\), and therefore includes a sink.  Put
\(\operatorname{Id}:\mathcal P(Q)\to\mathcal P(Q)\),
\(\operatorname{Id}(P)=P\).  For transformers
\(g,h:\mathcal P(Q)\to\mathcal P(Q)\), define their pointwise union and
intersection by
\[
 (g\cup h)(P):=g(P)\cup h(P),
 \qquad
 (g\cap h)(P):=g(P)\cap h(P).
\]
These are metalinguistic operations on predicate transformers, and every
union or intersection of transformers in Propositions~1--2 has this pointwise
standing.  For
\(h:\mathcal P(Q)\to\mathcal P(Q)\), put \(h^0=\operatorname{Id}\) and
\(h^{i+1}=h\circ h^i\).  Pointwise, the source defines
\[
 h^*(P)=\bigcup_{i\in\mathbb N}h^i(P),
 \qquad
 h^x(P)=\bigcap_{i\in\mathbb N}h^i(P).
\]
The superscript \(x\) is the source's literal name for the pointwise
intersection of all finite iterates, not multiplication and not an otherwise
undefined greatest-fixed-point binder.  No claim is made that the successive
sets form a descending chain for an arbitrary transformer.

Under the printed hypothesis that \(\to\) is image-finite---every state has
only finitely many immediate successors---Propositions~1--2 state:
\[
\begin{aligned}
\lvert\operatorname{ALL}(f)\rvert
 &=\widetilde{\operatorname{pre}}^{\,x}\lvert f\rvert,\\
\lvert\operatorname{SOME}(f)\rvert
 &=\bigl(\operatorname{Id}\cap
   (\operatorname{pre}\cup\widetilde{\operatorname{pre}})\bigr)^x
   \lvert f\rvert,\\
\lvert\operatorname{POT}(f)\rvert
 &=\operatorname{pre}^*\lvert f\rvert,\\
\lvert\operatorname{INEV}(f)\rvert
 &=\bigl(\operatorname{Id}\cup
   (\operatorname{pre}\cap\widetilde{\operatorname{pre}})\bigr)^*
   \lvert f\rvert.
\end{aligned}
\]
In the last line, the existential predecessor factor prevents a non-\(f\)
sink from entering vacuously.  The paper does not prove either proposition:
it says the fixed-point results are proved in detail in Sifakis's 1979 report.
Image-finiteness and this delegated-proof boundary are therefore parts of the
result \cite[pp.~347--348]{QueilleSifakis1982}.

\paragraph{What the analyzer actually computes.}
The analyzer assigns a Boolean variable to every IPN place, expresses
\(\operatorname{Init}\), \(\operatorname{after}a\), and
\(\operatorname{enable}a\) as predicates, reduces the net while protecting
places used by those predicates, computes \(\operatorname{pre}\), and
evaluates temporal subformulas in formula order.  The chapter does not print
the net-reduction rules, an IPN-to-\(\operatorname{pre}\) formula, or their
preservation proofs.

Its one complete calculation checks
\[
 \operatorname{Init}\Rightarrow
 \operatorname{ALL}\operatorname{POT}
 \operatorname{enable}(\mathsf{SENDER}).
\]
For every displayed Figure-2 place, a lower-case symbol
\(s_i,r_i,m_i,a_i\) is the Boolean characteristic predicate that the
like-named IPN place is marked.  In particular, the place predicate \(a_i\)
is not the action metavariable \(a\) used in the schemas above.  With
juxtaposition for Boolean conjunction, an overbar for Boolean negation, and
\(\cup\) for Boolean disjunction, the four translator-generated one-token
invariants used by the calculation are
\[
\begin{aligned}
s_1\bar s_2\bar s_3\cup\bar s_1s_2\bar s_3
 \cup\bar s_1\bar s_2s_3&=\mathsf{true},\\
r_1\bar r_2\bar r_3\cup\bar r_1r_2\bar r_3
 \cup\bar r_1\bar r_2r_3&=\mathsf{true},\\
m_1\bar m_2\cup\bar m_1m_2&=\mathsf{true},\\
a_1\bar a_2\cup\bar a_1a_2&=\mathsf{true}.
\end{aligned}
\]
In that notation,
\[
\begin{aligned}
\operatorname{Init}
 &=s_1m_1a_1r_1\bar Y\bar Z,\\
P_0=\operatorname{enable}(\mathsf{SENDER})
 &=s_1m_1\cup s_2m_1\cup s_2a_2\cup s_3,\\
\operatorname{pre}(P_0)
 &=s_3\cup r_2\cup r_3\cup a_2\cup m_2.
\end{aligned}
\]
Using these four one-token invariants, the source simplifies
\[
 P_1=P_0\cup\operatorname{pre}(P_0)=\mathsf{true}.
\]
Hence
\(\operatorname{POT}\operatorname{enable}(\mathsf{SENDER})\),
its \(\operatorname{ALL}\)-closure, and the whole implication all evaluate to
\(\mathsf{true}\).  None of the other property formulas is worked through
\cite[pp.~348--349]{QueilleSifakis1982}.

There is no asymptotic time or space bound and no theorem that the analyzer
always answers.  The source says there is no convergence-speed criterion:
the user chooses a maximum number of iterations, and if an iteration has not
finished by then the analyzer returns no answer.  The then-current
implementation supports Boolean, enumerated, and bounded-integer variables,
encoding the latter two into Booleans.  An iteration cap is therefore an
implementation boundary, not a proof by approximation.  This presentation
supplies a model-checking architecture and four predicate-transformer
identities at their stated image-finite strength; it does not supply CTL
satisfiability, a complete calculus, a general complexity theorem, or a CESAR
synthesis result \cite[pp.~348--350]{QueilleSifakis1982}.

\subsubsection{Emerson--Halpern's typed CTL* and expressiveness presentation}
\label{sec:p09-presentation-n-c-016-emerson-halpern}

The controlling source is the complete twenty-eight-page JACM article,
pp.~151--178.  An archived ACM publisher PDF was read page by page; its first
page carries ACM's conditional copying terms, not an open-content licence.
The POPL 1983 preliminary paper and the two 1984 reports named by the journal
are distinct source states.  This reconstruction therefore paraphrases the
1986 mathematical presentation and does not reproduce the scan, its figures,
or protected prose \cite[pp.~151--153]{EmersonHalpern1986}.

\paragraph{Why one untyped temporal syntax is not enough.}
The paper begins with Lamport's structure
\[
 M=(S,X,L),
\]
where \(S\neq\varnothing\) is a set of states,
\(X\neq\varnothing\) is a set of nonempty finite or infinite state
sequences, and \(L(s)\) is the set of atoms true at \(s\).  A path with
\(1+k\) states has length \(k\); an infinite path has length \(\omega\).
For a path \(x\), \(first(x)\) is its first state,
\(x^{succ}\) removes that state when possible and otherwise fixes a
length-zero path, and
\[
 x^0=x,\qquad x^{m+1}=(x^m)^{succ}.
\]
These operations define the suffixes used below.

Lamport's shared syntax has a primitive \(\Box p\), “always \(p\),”
and a different primitive \(\rightsquigarrow p\), “sometimes \(p\).”
The paper writes
\[
 \Diamond p:=\neg\Box\neg p
\]
for “not never \(p\).”  The distinction is necessary because the same
glyphs receive two semantic types.  Under the branching reading they are
state assertions:
\[
\begin{aligned}
 M,s\models_B\Box p
 &\Longleftrightarrow
 \forall x\in X\,
 \bigl(first(x)=s\Rightarrow
       \forall n\geq 0\;
       M,first(x^n)\models_Bp\bigr),\\
 M,s\models_B\rightsquigarrow p
 &\Longleftrightarrow
 \forall x\in X\,
 \bigl(first(x)=s\Rightarrow
       \exists n\geq 0\;
       M,first(x^n)\models_Bp\bigr).
\end{aligned}
\]
Under the linear reading they are assertions about one path:
\[
\begin{aligned}
 M,x\models_L\Box p
 &\Longleftrightarrow
   \forall n\geq0\;M,x^n\models_Lp,\\
 M,x\models_L\rightsquigarrow p
 &\Longleftrightarrow
   \exists n\geq0\;M,x^n\models_Lp.
\end{aligned}
\]
Thus \(\rightsquigarrow p\) and \(\Diamond p\) coincide on a single
linear path but not under the branching reading.

The first structure-class hypothesis used below is a property of the path set:
\(X\) is \emph{suffix closed} exactly when
\[
 \forall x\in X\;(x^{succ}\in X).
\]
It preserves each individual path's suffixes; unlike the two additional
closure properties defined below, it neither splices two paths nor takes an
infinite limit.  The distinction is used immediately in Theorems~1--2.

Definition~1 calls a formula model-valid when it is true at every state
under the branching reading or on every member of \(X\) under the linear
reading.  Definition~2 calls formulas under possibly different readings
\emph{strongly equivalent} over a structure class when their model-validity
agrees on every structure in the class.  Theorems~1--2 restate Lamport's
claims that branching \(\Diamond P\) has no strongly equivalent linear
assertion and that linear
\(\rightsquigarrow\Box P\) has no strongly equivalent branching assertion,
over suffix-closed structures.  Their proofs are not reproduced here.
Proposition~1 then exposes the coarseness of model-level comparison:
\(\rightsquigarrow P\land\neg P\) is satisfiable at an individual state or
path but is strongly equivalent to \(\mathsf{false}\), because no whole
model satisfies it everywhere \cite[pp.~153--155]{EmersonHalpern1986}.

\paragraph{The structure classes used by the comparison.}
Let \(S^*\) be the finite state strings, including the empty string,
\(S^+\) the nonempty finite state strings, and \(S^\omega\) the infinite state
strings.  Juxtaposition denotes string concatenation.  In a fusion
decomposition \(x_1sy_1\), take
\(x_1\in S^*\), \(s\in S\), and \(y_1\in S^*\cup S^\omega\), and likewise
for the second decomposition, whenever the displayed concatenations are
well-formed nonempty paths.  In a limit decomposition, take each
\(x_i\in S^+\) and each \(y_i\in S^*\cup S^\omega\), again only when every
displayed concatenation is a well-formed path.  The following implications are
universally quantified over these compatible decompositions and define
properties of \(X\), not temporal operators:
\[
\begin{array}{ll}
\text{suffix closed}
 &x\in X\Rightarrow x^{succ}\in X,\\[2pt]
\text{fusion closed}
 &x_1sy_1,x_2sy_2\in X\Rightarrow x_1sy_2\in X,\\[2pt]
\text{limit closed}
 &\bigl(\forall n\geq1\;x_1\cdots x_ny_n\in X\bigr)
   \Rightarrow x_1x_2x_3\cdots\in X.
\end{array}
\]
For an infinite-path set \(X\subseteq S^\omega\), the set \(X\) is
\(R\)-generable when some total relation \(R\subseteq S\times S\) satisfies
\[
 X=\{x\in S^\omega:\forall n\in\mathbb N\;((x(n),x(n+1))\in R)\}.
\]
A structure is \emph{state complete} when every \(s\in S\) begins a path
in \(X\).  For \(X\subseteq S^\omega\), the characterization is
\[
 \text{\(R\)-generable}
 \quad\Longleftrightarrow\quad
 \text{suffix closed, fusion closed, limit closed, and state complete}.
\]
For the forward implication, suffixes and fusions of \(R\)-paths are
\(R\)-paths because each adjacent pair remains an \(R\)-edge.
Every adjacent pair of a limiting path already occurs in one of its
extendible finite prefixes, so that path is also an \(R\)-path.
Totality supplies an infinite \(R\)-path from each state, using dependent
choice in the ambient metatheory.

For the converse, define
\[
 R=\{(s,t)\in S^2:\text{some }x\in X
                         \text{ begins with }st\}.
\]
State completeness makes \(R\) total, since each path in \(X\) is infinite.
Suffix closure shows that every adjacent pair of every member of \(X\)
belongs to \(R\).  Thus \(X\) is contained in the set of all \(R\)-paths.
Let \(s_0s_1\cdots\) be an arbitrary \(R\)-path.
Its first edge extends to a member of \(X\).
If \(s_0\cdots s_n\) extends to a member of \(X\), fuse that extension at
\(s_n\) with a witness for \((s_n,s_{n+1})\in R\).
This extends \(s_0\cdots s_{n+1}\) to a member of \(X\).
Induction therefore makes every nonempty finite prefix extendible;
limit closure puts the entire path in \(X\).
This proves the reverse containment and the characterization.

State completeness cannot be omitted: adjoining a state that appears in
no member of \(X\) preserves the three closure properties but prevents
generation by a total relation on the enlarged \(S\).
Fusion closure supports splicing futures at a shared state, limit closure
supports an infinite limiting computation, and state completeness supplies
paths at every state.
Lamport's original
Theorem~1 witnesses need not be fusion closed and his Theorem~2 witnesses
need not be limit closed.  Section~4 therefore re-establishes the
separations over \(R\)-generable structures
\cite[pp.~155--156]{EmersonHalpern1986}.

\paragraph{The two formula types of CTL*.}
The source now replaces the shared syntax by mutually defined state and path
formulas:
\[
\begin{aligned}
 \varphi &::=
 P\mid\neg\varphi\mid(\varphi\land\varphi)
 \mid A\psi\mid E\psi,\\
 \psi &::=
 \varphi\mid\neg\psi\mid(\psi\land\psi)
 \mid F\varphi\mid F\psi
 \mid X\varphi\mid X\psi\\
 &\qquad
 \mid(\varphi U\varphi)\mid(\psi U\psi)
 \mid\widetilde F\varphi\mid\widetilde F\psi.
\end{aligned}
\]
Here \(\varphi\) has standing as a state assertion and \(\psi\) as a path
assertion; \(A/E\) are the only constructors that turn a path assertion into
a state assertion.  This type distinction is used by every sublanguage and
nondefinability result below.  It differs from Lamport's syntax exactly by
making the state/path standing visible, rather than asking one formula to
change type with its interpretation.

For \(M=(S,X,L)\),
\[
\begin{aligned}
 s\models A\psi
 &\Longleftrightarrow
   \forall x\in X\,
   (first(x)=s\Rightarrow x\models\psi),\\
 s\models E\psi
 &\Longleftrightarrow
   \exists x\in X\,
   (first(x)=s\land x\models\psi).
\end{aligned}
\]
A state formula used as a path formula is evaluated at the path's first
state.  The temporal clauses are
\[
\begin{aligned}
 x\models F\psi
 &\Longleftrightarrow
   \exists i\geq0\;x^i\models\psi,\\
 x\models X\psi
 &\Longleftrightarrow
   |x|\geq1\land x^1\models\psi,\\
 x\models\psi_1U\psi_2
 &\Longleftrightarrow
   \exists i\geq0\,
   \bigl(x^i\models\psi_2
     \land\forall j<i\;x^j\models\psi_1\bigr),\\
 x\models\widetilde F\psi
 &\Longleftrightarrow
   \text{\(x^i\models\psi\) for infinitely many distinct \(i\)}.
\end{aligned}
\]
The duplicated source clauses for state operands replace
\(x^i\models\varphi\) by \(first(x^i)\models\varphi\).
The \(F\) and strong-\(U\) clauses include the current position, while \(X\)
is false at the last state of a finite path.  A footnote makes a further
distinction: path formulas are semantically evaluated on all nonempty finite
or infinite state sequences over \(S\), paths in the preceding sense; only
\(A/E\) restrict their quantified paths to members of \(X\).

The source defines
\[
\begin{gathered}
 G\psi:=\neg F\neg\psi,\qquad
 \widetilde G\psi:=\neg\widetilde F\neg\psi,\\
 A\psi:=\neg E\neg\psi,\qquad
 F\psi:=(\mathsf{true}\ U\psi),\qquad
 \widetilde F\psi:=GF\psi.
\end{gathered}
\]
Thus \(E,X,U,\neg,\land\) can serve as primitives.  The full state language
is CTL*.  The exact named fragments are
\[
\begin{gathered}
 L(F),\quad L(F,U),\quad L(F,X,U),\\
 B(F),\quad B(F,X,U),\quad
 B(F,X,U,\land,\neg),\quad
 B(F,X,U,\widetilde F,\land,\neg).
\end{gathered}
\]
The \(L\)-languages contain path formulas without state-level path
quantification.  \(B(F)\) permits \(A/E\) followed by \(F\) applied directly
to state formulas; \(B(F,X,U)\) adds the corresponding direct \(X/U\)
forms.  Adding \(\land,\neg\) in the fragment name permits Boolean
combinations \emph{inside} path formulas, not merely the state-level Boolean
operations already present.  The source identifies \(B(F,X,U)\) with CTL
and defines \(|p|\) as the number of symbols in the formula string
\cite[pp.~156--158]{EmersonHalpern1986}.

\paragraph{From whole-model validity to formula equivalence.}
For a Lamport formula \(p\), let \(p^L\) replace
\(\Box,\rightsquigarrow\) by \(G,F\), and let \(p^B\) replace them by
\(AG,AF\).  Proposition~2 gives the two pointwise satisfaction
equivalences.  In particular, the linear principle
\[
 FP\lor G\neg P
\]
is valid, whereas \(AFP\lor AG\neg P\) is not.  Its correctly typed
branching assertion is \(A(FP\lor G\neg P)\).

Definition~3 associates with a path language \(L\) the state language
\[
 B(L):=\{A\psi:\psi\in L\}.
\]
This is a metalinguistic language construction, not an object-language
operator, and the source warns that it is not closed under semantic negation
or disjunction.  Definition~4 then compares state formulas pointwise at
every state of every structure in the chosen class.

Proposition~3 proves that any path formula \(p\), considered model-wide, is
strongly equivalent to \(Ap\).  For a structure class \(\mathcal C\), write
\(p\equiv_s^{\mathcal C}q\) when \(p\) and \(q\) have the same model-validity
status in every member of \(\mathcal C\), and write
\(p\equiv^{\mathcal C}q\) when the two state formulas agree at every state of
every member of \(\mathcal C\).  Let \(\mathcal F\) be the class of
fusion-closed, state-complete structures.  Proposition~4 proves
\[
 p\equiv_s^{\mathcal F}q
 \quad\Longleftrightarrow\quad
 AGp\equiv^{\mathcal F}AGq
\]
for state formulas \(p,q\).  Its forward proof restricts a fusion-closed
model to the states occurring on paths from the selected state; its reverse
proof uses state completeness.  The source demonstrates nonredundancy of
the two hypotheses with
\[
\begin{array}{lll}
\text{no fusion closure:}
 &p=P\land EFEX\neg P,&q=\mathsf{false},\\
\text{no state completeness:}
 &p=AGEF\mathsf{true}\land EF\mathsf{true},&
 q=p\lor AG\mathsf{false}.
\end{array}
\]
Corollary~1 gives the mixed path/state form
\[
 p\equiv_s^{\mathcal F}q
 \quad\Longleftrightarrow\quad
 AGAp\equiv^{\mathcal F}AGq.
\]
Definition~5 finally orders languages by the existence of pointwise
equivalent state formulas \cite[pp.~158--160]{EmersonHalpern1986}.

\paragraph{The exact expressiveness map.}
All comparisons in Section~4 are over \(R\)-generable structures.  The
paper says its proofs also work over suffix-closed, fusion-closed structures,
but leaves that extension to the reader.  Figure~1 is equivalently the three
chains and one cross-edge
\[
\begin{aligned}
 B(L(F))
 &<B(L(F,U))
 <B(L(F,X,U))
 <CTL^*,\\
 B(F)
 &<B(F,X,U)
 =B(F,X,U,\land,\neg)\\
 &<B(F,X,U,\widetilde F,\land,\neg)
 <CTL^*,\\
 B(L(F))
 &<B(F,X,U,\widetilde F,\land,\neg).
\end{aligned}
\]
Any two displayed languages not connected by the transitive closure of
these \(<\) and \(=\) relations are incomparable.  The absence of another
edge is therefore mathematical content, not blank space in the diagram
\cite[pp.~160--163]{EmersonHalpern1986}.

Theorem~3 supplies one complete separating proof.  Let \(M\) have states
\(s_1,s_2\),
\[
 R=\{(s_1,s_1),(s_1,s_2),(s_2,s_2)\},
 \qquad s_1\models\neg P,\quad s_2\models P,
\]
and let \(M'\) be its one-state restriction to \(s_1\).
Then \(M,s_1\models EFP\) but \(M',s_1\not\models EFP\).
If \(EFP\) were equivalent to \(Aq\), truth of \(Aq\) in \(M\) would survive
deletion of paths and hold in \(M'\), a contradiction.  Hence \(EFP\) is
not expressible by any \(B(L(F,X,U))\) formula.  The invariant is preservation
of a universal path assertion under path deletion; existential reachability
does not have that preservation property.

Theorem~4 proves
\[
 B(L(F))\leq
 B(F,X,U,\widetilde F,\land,\neg).
\]
The Appendix first normalizes every linear \(F/G\) formula into proof-local
forms and their finite conjunctions and disjunctions.  Those auxiliary forms
are not the associated-language constructor \(B(L)\).  The normalization uses
propositional formulas and, for positive arity \(n\geq1\), the segment forms
\[
 [p_1]:=Gp_1,
 \qquad
 [p_1,\ldots,p_n]:=p_1U[p_2,\ldots,p_n]\quad(n\geq2),
\]
\(GFp\), and the source's nested \(F\)-forms, and then translates their Boolean
combinations.  Conjoined segment forms are replaced by a
disjunction over consistent interleavings.  The final induction decreases
the number of nested \(F\)-subformulas by choosing which eventuality is met
first and in which segment.  No \(X\) occurs, so the proof
lands in the displayed rich branching language even after every
\(X\)-forming clause is removed
\cite[Theorem~4, pp.~167--169]{EmersonHalpern1986}.

\paragraph{The separating model families.}
Theorem~5 states that
\[
 A[F(P\land XP)]
\]
is not equivalent to any
\(B(F,X,U,\widetilde F,\land,\neg)\) formula.
Figures~2--3 start with
\(M_1\), a \(P,P,\neg P\) chain ending in a self-loop, and \(N_1\), a
\(P,\neg P\) chain ending in a self-loop.
The unique path of \(M_1\) witnesses \(F(P\land XP)\) at its first
position.  The unique path of \(N_1\) has no consecutive \(P\)-states.

Here is the recursive construction in Figure~3.  For each new model,
take a new \(P\)-root \(a_{i+1}\), a new \(\neg P\)-state \(b_{i+1}\),
and disjoint copies \(M_i'\), \(N_i'\), and one further lower-level model.
Add edges from \(a_{i+1}\) to \(b_{i+1}\) and to the roots of both
\(M_i'\) and \(N_i'\).  In \(M_{i+1}\), the only successor of
\(b_{i+1}\) is the root of the further copy \(M_i\); in \(N_{i+1}\),
it is the root of the further copy \(N_i\).  Retain all edges and
labels within the copies, and add no other edges.

Every path taking one of the two direct copy edges begins with \(P,P\).
Every remaining path in \(M_{i+1}\) reaches \(M_i\), where the induction
hypothesis supplies a consecutive \(P\)-pair.  In \(N_{i+1}\), follow
\(b_{i+1}\) into \(N_i\) and then a path without such a pair.
The newly traversed labels \(P,\neg P,P\) add no consecutive pair.
This proves the first of the three invariants:
\[
\begin{array}{ll}
(1)&M_i,a_i\models A[F(P\land XP)],
     \quad
     N_i,a_i\not\models A[F(P\land XP)];\\
(2)&\text{on these families, every }
 B(F,X,U,\widetilde F,\land,\neg)\text{ formula}\\
&\quad\text{has an equivalent }B(F,X,U)\text{ formula};\\
(3)&|p|\leq i\Rightarrow
 (M_i,a_i\models p\Longleftrightarrow N_i,a_i\models p)
 \quad\text{for }p\in B(F,X,U).
\end{array}
\]
For invariant~(3), the Appendix prints the \(EXq,E[qUr],A[qUr]\) induction
cases.  Atoms agree because the paired roots have the same labels.
Conjunction and negation preserve agreement by the induction hypotheses
for their strictly smaller operands.
For invariant~(2), every path in every model eventually stays at a
\(\neg P\)-self-loop.  These terminal loops have identical future
structures.  Hence, on these models, each
\(\widetilde F/\widetilde G\) subformula has the constant truth value
determined by its operand on that terminal loop.  Eliminating these
operators from the inside outward leaves Boolean path combinations;
apply
Theorem~8(a)'s Boolean-path elimination,
\[
 B(F,X,U)=B(F,X,U,\land,\neg).
\]
If a formula in the rich branching fragment defined the witness,
invariant~(2) would give an equivalent \(B(F,X,U)\) formula on these
families.  Choosing \(i\) at least its size would make that formula agree
at the two roots by invariant~(3), contrary to invariant~(1).

Theorem~6 gives a second formula not definable in
\(B(F,X,U,\widetilde F,\land,\neg)\):
\[
 \Theta:=
 E[((P_1UP_2)\lor(Q_1UQ_2))UR].
\]
Figures~4--5 begin with two five-state chains whose middle labels order
\(P_2,Q_1\) differently.  At the next stage, both structures receive an
identically labelled \(P_1,P_2,Q_1,Q_2\) prefix, but copies of the preceding
\(M_i\) and \(N_i\) are attached in different arrangements.  The source
records \(M_i,a_i\models\Theta\), \(N_i,a_i\not\models\Theta\), and the same
reduction and bounded-formula invariants as for Theorem~5.  It leaves the
repeated induction to the reader.  The same bounded-size contradiction
establishes the stated nondefinability.

The dual
\(\neg\Theta\) is in \(B(L(F,X,U))\).  If it were expressible in the rich
branching fragment, closure of that fragment under state negation would make
\(\Theta\) expressible there, contradicting Theorem~6.  Hence
\(B(L(F,X,U))\) is not contained in the rich branching fragment
\cite[Theorems~5--6, pp.~161--163, 169--171]{EmersonHalpern1986}.

Theorem~7 uses Figures~6--8 to show
\[
 E\widetilde FP
\]
is not expressible in \(B(F,X,U,\land,\neg)\).  Its \(M_i\) and \(N_i\)
families use \(P\)-states \(a_i,c_i\) and \(\neg P\)-states \(b_i,d_i\).
Each \(M_{i+1}\) adds a return edge from \(b_{i+1}\) to \(a_{i+1}\) before
descending into \(N_i\); \(N_{i+1}\) has the corresponding descending
prefix without that return edge.  Hence
\(M_i,a_i\models E\widetilde FP\) and
\(N_i,c_i\not\models E\widetilde FP\).  The first induction proves that,
beyond \(|p|\), truth of a \(B(F,X,U)\) formula at all \(c_j\) states is
equivalent to truth at some such state, and likewise for the \(d_j\).  A
second induction proves
\[
 i\geq|p|\Rightarrow
 \bigl(a_i\models p\Longleftrightarrow c_i\models p\bigr)
 \land
 \bigl(b_i\models p\Longleftrightarrow d_i\models p\bigr).
\]
Theorem~8(a)'s cited Boolean-path elimination transfers the invariant to
\(B(F,X,U,\land,\neg)\), while \(E\widetilde FP\) separates the roots.
The Appendix prints the \(EX,E[U],A[U]\) cases and leaves Boolean cases to
the reader.

Theorem~8, cited from the authors' earlier paper, states
\[
 B(F,X,U)=B(F,X,U,\land,\neg)
\]
and says \(E(PUQ)\) is not expressible in \(B(F)\).  Theorem~9, cited from
Process Logic, says \(PUQ\) is not expressible in \(L(F)\) and \(XP\) is
not expressible in \(L(F,U)\).  This article supplies neither cited proof
nor its separating structures.  Those citations, rather than new local
proofs, establish the strict edges from \(B(F)\), \(B(L(F))\), and
\(B(L(F,U))\) in this source presentation
\cite[pp.~161--175]{EmersonHalpern1986}.

\paragraph{What the comparison says about specifications.}
For the transition under discussion, \(\mathit{enabled}\) and
\(\mathit{executed}\) are atomic state predicates: the first holds when the
transition is enabled, and the second when it executes.  Lamport's
strong-eventual-fairness assertion is represented by
\[
 A(\widetilde G\neg\mathit{enabled}\lor F\mathit{executed}).
\]
Because \(\widetilde G\neg\mathit{enabled}\) says that enabledness occurs only
finitely often, the formula says that infinite recurring enabledness entails
eventual execution.  It is a specification formula, not a fairness condition
built into \(X\).

\paragraph{The PL translation and its proof boundary.}
To a CTL* structure \(M=(S,X,L)\), the paper associates a PL path model in
which one atomic program \(A\) denotes exactly \(X\) and the atomic labels
agree.  The operator \(\mathsf f\) tests the unique length-zero initial prefix,
while \(\mathrm{suf}\) selects a proper suffix after an interval satisfying
its left operand.  Taking \(E,X,U\) as primitive, the source prints
\[
\begin{aligned}
 P'&=P,
 &( \neg p)'&=\neg p',
 &(p\land q)'&=p'\land q',\\
 (pUq)'&=q'\lor(p'\mathbin{\mathrm{suf}}q'),
 &(Ep)'&=\mathsf f(\langle A\rangle p'),
 &(Xp)'&=\mathsf{false}\mathbin{\mathrm{suf}}p'.
\end{aligned}
\]
Proposition~5 asserts preservation for all path formulas on
\(x\in X\) and all state formulas at \(s\in S\).  The paper calls the proof
a straightforward structural induction but does not print its cases.  Its
Appendix supplies only the needed test-free-PDL and path-model definitions,
not a complete new proof theory for PL
\cite[pp.~164, 175--176]{EmersonHalpern1986}.

\paragraph{The MPL semantics and translation.}
\label{sec:p09-mpl-translation}
MPL formulas are generated from atoms by \(\neg,\land,X,U,\Diamond\).
An MPL formula is evaluated at \(M,x,y\), where \(x\in X\) and the
\emph{stage} \(y\) is a nonempty finite prefix of \(x\).  Prefix order is
\(y\leq z\).  The term is needed because MPL may change to another path
sharing the already traversed prefix; it stands for that finite prefix, not
for a state or a new temporal connective, and Lemmas~1--2 use it to relate
arbitrary structures to suffix- and fusion-closed ones.

An atom \(P\) is true at \(M,x,y\) exactly when \(P\in L(last(y))\),
where \(last(y)\) is the final state of the finite string \(y\).
Negation and conjunction have their ordinary Boolean clauses.
In the following temporal and path-choice clauses, \(z,w\) range over
nonempty finite prefixes of \(x\):
\[
\begin{aligned}
 M,x,y\models pUq
 \Longleftrightarrow{}&
 \exists z\bigl(
 y\leq z\leq x\land M,x,z\models q\\
 &\qquad\land
 \forall w\,
 (y\leq w<z\Rightarrow M,x,w\models p)\bigr),\\
 M,x,y\models Xp
 \Longleftrightarrow{}&
 \exists z\bigl(
 M,x,z\models p\land y<z\leq x\\
 &\qquad\land
 \neg\exists w\,(y<w<z)\bigr),\\
 M,x,y\models\Diamond p
 \Longleftrightarrow{}&
 \exists x'\in X\,
 (y\leq x'\land M,x',y\models p).
\end{aligned}
\]
Thus \(pUq\) requires \(p\) strictly before the witnessing \(q\)-stage,
and \(Xp\) requires the next stage.  A current \(q\)-stage witnesses
\(pUq\) regardless of the current value of \(p\).  At the last stage
of a finite path, \(Xp\) is false.

The source motivates its closure hypotheses with two further exact
separations.  The CTL* formula
\[
 EGX\mathsf{true}\land\neg EXEGX\mathsf{true}
\]
is satisfiable only outside suffix-closed structures, and
\[
 EFEFp\land\neg EFp
\]
is satisfiable only outside fusion-closed structures.  By contrast,
Lemma~1 says that every pointed MPL model has a satisfaction-equivalent point
in a structure possessing both closure properties.  These claims explain why
Theorem~10 restricts its embedding.

Lemma~1's Appendix construction unwinds the paths extending a stage into a
fresh-state tree and closes it under suffix; for an uncountable path family
the authors say they appear to need the well-ordering principle.

The translation on a suffix-closed, fusion-closed structure is as follows.
Use the CTL* primitives \(E,X,U,\neg,\land\), and define \(\tau\)
recursively by
\[
\begin{aligned}
 \tau(P)&=P,&
 \tau(\neg\alpha)&=\neg\tau(\alpha),&
 \tau(\alpha\land\beta)&=\tau(\alpha)\land\tau(\beta),\\
 \tau(X\alpha)&=X\tau(\alpha),&
 \tau(\alpha U\beta)&=\tau(\alpha)U\tau(\beta),&
 \tau(E\alpha)&=\Diamond\tau(\alpha).
\end{aligned}
\]
A state formula embedded in a path formula has the same translation.
The other operators are expanded by their stated definitions.

For \(x=s_0s_1\cdots\in X\), write \(y_i=s_0\cdots s_i\) at each
actual position \(i\) of \(x\).  The preservation statements are
\[
\begin{aligned}
 M,x,y_i\models_{\mathrm{MPL}}\tau(\psi)
 &\Longleftrightarrow M,x^i\models_{\mathrm{CTL}^*}\psi
 &&\text{for path formulas }\psi,\\
 M,x,y_i\models_{\mathrm{MPL}}\tau(\varphi)
 &\Longleftrightarrow M,s_i\models_{\mathrm{CTL}^*}\varphi
 &&\text{for state formulas }\varphi.
\end{aligned}
\]
We prove them simultaneously by structural induction.
The atomic case is the definition of \(L(s_i)\).
Negation and conjunction follow from the induction hypotheses, and
embedding a state formula into the path language uses
\(first(x^i)=s_i\).

For \(X\psi\), a proper prefix successor of \(y_i\) exists exactly
when \(i+1\) is an actual position of \(x\); it is then \(y_{i+1}\).
The induction hypothesis at that stage gives the CTL* next-time clause.
If no such position exists, both sides are false.
For \(\psi U\chi\), the MPL witness is a prefix \(y_j\) with
\(j\geq i\).  The prefixes \(w\) with \(y_i\leq w<y_j\) are
exactly \(y_i,\ldots,y_{j-1}\).  Applying the induction hypotheses at
these finitely many stages gives the CTL* until clause on \(x^i\).
Conversely, a CTL* until witness on a finite path can always be chosen
at an actual position: if its index is past the end, replace it by
the final index.  The suffix there is the same, and the requirements
before it only decrease.  It therefore gives an MPL witness \(y_j\).

For \(E\psi\), an MPL witness is \(x'\in X\) extending \(y_i\).
Suffix closure gives \(v=(x')^i\in X\) with \(first(v)=s_i\).
The induction hypothesis on \(x'\) at \(y_i\) turns its MPL truth
into \(M,v\models_{\mathrm{CTL}^*}\psi\), witnessing the state formula.
Conversely, let \(v\in X\) start at \(s_i\) and satisfy \(\psi\).
Fusion closure, applied to \(x\) at \(s_i\) and to \(v\) at its first
state, gives
\[
 x'=s_0\cdots s_{i-1}v\in X.
\]
For \(i=0\), the displayed prefix before \(v\) is empty.
The path \(x'\) extends \(y_i\) and \((x')^i=v\), so the induction
hypothesis supplies the MPL witness.  This proves both directions of
the path-choice step and completes the induction.

The same calculation also proves Lemma~2's suffix invariance.
Every MPL formula \(p\) is the translation of a CTL* path formula,
obtained recursively by replacing \(\Diamond\) with \(E\).
Apply preservation at \(y_i\) on \(x\) and at the first stage on
\(x^i\); both sides equal the CTL* truth of that path formula on \(x^i\).
Thus, writing \(x/y_i=x^i\), satisfaction at \(M,x,y_i\) agrees
with satisfaction at the first stage of \(M,x/y_i\).
In particular the displayed translation proves Theorem~10 at every
path's first stage and at every state that starts a path in \(X\)
\cite[Lemma~2 and Theorem~10, pp.~165--166]{EmersonHalpern1986}.

The final comparison says that Lehmann--Shelah's TC replaces \(A\) by a
measure-one path operator.  It reports that TC and MPL have the same complete
axiomatizations and hence the same valid and satisfiable formulas, then says
one would \emph{expect} a CTL*/MPL-to-TC translation.  No such translation
is constructed in this paper
\cite[pp.~165--166, 176--177]{EmersonHalpern1986}.

\paragraph{The terminal-server example and the scope of the conclusion.}
For a terminal-server process \(P_i\), let \(\mathit{inNCS}_i\) and
\(\mathit{inTRY}_i\) be atomic state predicates holding exactly when
\(P_i\) is respectively in its noncritical or trying region.  The article
requires a possible path that remains forever in \(\mathit{inNCS}_i\), another
possible path that reaches \(\mathit{inTRY}_i\), and a universal constraint
that every path either remains or eventually tries:
\[
 EG(\mathit{inNCS}_i)\land EF(\mathit{inTRY}_i)\land
 A(G(\mathit{inNCS}_i)\lor F(\mathit{inTRY}_i)).
\]
The linear candidate
\(A(G(\mathit{inNCS}_i)\lor F(\mathit{inTRY}_i))\) admits a degenerate
model in which every path eventually reaches \(\mathit{inTRY}_i\) and no path
remains in \(\mathit{inNCS}_i\).  Explicit existential path quantification
therefore has an actual use in a nondeterministic specification and in excluding a
single-path synthesized implementation.

Section~6 nevertheless presents an application argument, not a dominance
theorem.  It says linear time is often adequate and simpler for manual
verification of an existing program when all executions are quantified
universally; it favors branching constructs for model checking, negating
universal path claims, and requiring alternative behaviors.  The paper
concentrates on expressiveness and refers elsewhere for satisfiability and
finite-model truth complexity.  It supplies no CTL state-labelling
algorithm, no general complexity bound, and no synthesis correctness theorem
\cite[pp.~166--167]{EmersonHalpern1986}.

\subsubsection{Clarke--Grumberg--Peled's 1999 acquisition boundary}
\label{sec:p09-presentation-n-c-016-clarke-grumberg-peled}

The audited official routes did not yield an anonymously readable complete
first-edition body or a public author-accepted manuscript as of 22 August
2026.  This is not a claim that the book is unavailable: physical copies,
access-restricted Internet Archive lending, and institution-only EBSCO access
exist.  Google Books is a searchable partial preview, HathiTrust's
standard-identifier queries returned no item, and the MIT Press first-edition
page is an out-of-print product record without a body file.  Its contents page
locates the starts of CTL/CTL* at p.~27, explicit model checking at p.~35,
BDDs at p.~51, and symbolic model checking at p.~61, but locator metadata
supports no literal algorithm, proof, example, or complexity claim
\cite{ClarkeGrumbergPeled1999}.

The complete 156-page Clarke--Schlingloff 2001 handbook chapter, the
Clarke--Grumberg--Long 1996 survey, the expanded five-author 2018 second
edition, and the Baier--Katoen 2008 book presented next are different works or
source states; none supplies the missing first-edition text.  Accordingly,
the four complexity classifications in
Theorem~\ref{thm:p09-ctl-star-complexity} retain their independent citations;
none is attributed to the unacquired 1999 body at source-content strength.

\subsubsection{Baier--Katoen's CTL labelling presentation}
\label{sec:p09-presentation-n-c-016-baier-katoen}

This presentation follows the complete Baier--Katoen book.  Its copyright leaf
says \emph{All rights reserved}; the pages and protected expression are not
reproduced here.  The relevant page map also needs care:
Chapter~6 begins on printed p.~313, CTL syntax begins on p.~317, the model-
checking section occupies pp.~341--358, and the requested CTL* section begins
on p.~422.  Printed pp.~311--312 are Chapter~5 exercises, and pp.~417--421 are
still symbolic CTL rather than CTL* \cite{BaierKatoen2008}.

\paragraph{Inherited transition systems and paths.}
The book does not replace its earlier transition-system definition by an
unlabelled Kripke tuple.  It retains
\[
 TS=(S,Act,{\to},I,AP,L),
\]
where \({\to}\subseteq S\times Act\times S\), \(I\subseteq S\), and
\(L:S\to2^{AP}\).  It calls the system finite when \(S\), \(Act\), and
\(AP\) are finite.  For \(s\in S\), put
\[
 \operatorname{Post}(s,\alpha)
   =\{t\in S:s\xrightarrow{\alpha}t\},
 \qquad
 \operatorname{Post}(s)
   =\bigcup_{\alpha\in Act}\operatorname{Post}(s,\alpha).
\]
A terminal state has \(\operatorname{Post}(s)=\varnothing\).  Chapter~6's
standing hypothesis is literally that the transition system has no terminal
states:
\[
 \forall s\in S\;\bigl(\operatorname{Post}(s)\ne\varnothing\bigr).
\]
The source does not introduce the term \emph{total transition system} as a new
Chapter~6 term, although this displayed condition is the corresponding graph
property.

The source tuple and the shared handbook method have different types, so an
explicit bridge is necessary.  Forget the action name, but not the transition,
by defining
\[
 sR_0t\iff\exists\alpha\in Act\;(s\xrightarrow{\alpha}t),
 \qquad
 V_0(a)=\{s\in S:a\in L(s)\}\quad(a\in AP).
\]
These are handbook-supplied metalinguistic definitions: \(R_0\) is a binary
relation among the source states and \(V_0:AP\to\mathcal P(S)\) is its
propositional valuation; neither is a new Baier--Katoen transition system or
source term.  For this instantiation, take the host proposition alphabet
\(\mathsf{Prop}\) to be exactly \(AP\), so \(V_0\) has the host's required
type.

A Baier--Katoen CTL state formula \(\Phi\) is built from \(true\), atoms
\(a\in AP\), negation, conjunction, and the four immediate path-quantified
next/until forms.  Define the metalinguistic translation \(\operatorname{tr}\)
into the
host CTL syntax recursively by
\[
\begin{aligned}
 \operatorname{tr}(true)&:=\top,&
 \operatorname{tr}(a)&:=a,\\
 \operatorname{tr}(\neg\Phi)&:=\neg\operatorname{tr}(\Phi),&
 \operatorname{tr}(\Phi\land\Psi)
   &:=\operatorname{tr}(\Phi)\land\operatorname{tr}(\Psi),\\
 \operatorname{tr}(\exists\bigcirc\Phi)
   &:=\mathsf{EX}\operatorname{tr}(\Phi),&
 \operatorname{tr}(\forall\bigcirc\Phi)
   &:=\mathsf{AX}\operatorname{tr}(\Phi),\\
 \operatorname{tr}(\exists(\Phi\mathbin{U}\Psi))
   &:=\mathsf E[\operatorname{tr}(\Phi)\mathbin{\mathsf U}
                   \operatorname{tr}(\Psi)],&
 \operatorname{tr}(\forall(\Phi\mathbin{U}\Psi))
   &:=\mathsf A[\operatorname{tr}(\Phi)\mathbin{\mathsf U}
                   \operatorname{tr}(\Psi)].
\end{aligned}
\]
This translation is necessary because the source and host use different
constructor glyphs; it changes no atom, Boolean structure, or temporal
meaning.  It is not an object-language operator.  When \(S\ne\varnothing\)
and the source transition system is finite, together with
\(\mathcal K_0=(S,R_0,V_0)\), it supplies the exact typed input to
Method~\ref{def:p09-ctl-labelling}.
Directly from the definitions,
\[
 \operatorname{Post}(s)=\{t:sR_0t\},
 \qquad
 \operatorname{Pre}^{R_0}_{\exists}(X)
   =\{s:\operatorname{Post}(s)\cap X\ne\varnothing\}.
\]
If \(S\ne\varnothing\), the no-terminal-state condition makes \(R_0\) total;
when the source transition system is finite,
\(\mathcal K_0=(S,R_0,V_0)\) is therefore a finite total transition structure
of the host method's type.  If \(S=\varnothing\), then \(I=\varnothing\), every
satisfaction set is empty, and system satisfaction is vacuous; the host method
is not instantiated in that case.
If \(K\) counts the source's action-labelled transition triples, then
\(\lvert R_0\rvert\le K\): distinct action-labelled transitions can collapse
to the same binary edge.  Thus the shared theorem's binary-edge parameter and
Theorem~6.30's source-transition parameter must not be identified literally.

A finite path fragment is a sequence \(s_0\ldots s_n\) with
\(s_i\in\operatorname{Post}(s_{i-1})\) for \(0<i\le n\); the infinite case
is defined analogously.  For an infinite path
\(\pi=s_0s_1\ldots\), write
\(\pi[i..]=s_is_{i+1}\ldots\) for its suffix from position \(i\).  A maximal
path fragment is infinite or ends in a terminal state.  The set
\(\operatorname{Paths}(s)\) contains the maximal
fragments beginning at \(s\), whether or not \(s\in I\).  A path of the
transition system, by contrast, is initial as well as maximal.  Under the
no-terminal-state hypothesis, every member of \(\operatorname{Paths}(s)\) is
infinite \cite[defs.~2.1, 2.3--2.4 and 3.4--3.6; pp.~20, 23,
95--96]{BaierKatoen2008}.

\paragraph{Typed CTL syntax and semantics.}
For fixed \(AP\), Definition~6.1 gives the state-formula and path-formula
grammars
\[
\begin{aligned}
 \Phi&::=true\mid a\mid\Phi_1\land\Phi_2\mid\neg\Phi
          \mid\exists\varphi\mid\forall\varphi,
          &&a\in AP,\\
 \varphi&::=\bigcirc\Phi\mid\Phi_1\mathbin{U}\Phi_2.
\end{aligned}
\]
Thus a CTL temporal operator occurs immediately below a path quantifier; CTL
does not permit arbitrary Boolean or temporal nesting at path-formula level.
The displayed derived operators include
\[
\begin{aligned}
 \exists\Diamond\Phi&:=\exists(true\mathbin{U}\Phi),&
 \forall\Diamond\Phi&:=\forall(true\mathbin{U}\Phi),\\
 \exists\Box\Phi&:=\neg\forall\Diamond\neg\Phi,&
 \forall\Box\Phi&:=\neg\exists\Diamond\neg\Phi.
\end{aligned}
\]

For \(\pi=s_0s_1\ldots\in\operatorname{Paths}(s)\), Definitions~6.4--6.5
give
\[
\begin{aligned}
 s\models\exists\varphi
   &\iff \text{some }\pi\in\operatorname{Paths}(s)\text{ satisfies }\varphi,\\
 s\models\forall\varphi
   &\iff \text{every }\pi\in\operatorname{Paths}(s)\text{ satisfies }\varphi,\\
 \pi\models\bigcirc\Phi
   &\iff s_1\models\Phi,\\
 \pi\models\Phi\mathbin{U}\Psi
   &\iff \exists j\ge0\;\bigl(s_j\models\Psi\ \land\
      \forall k<j\;(s_k\models\Phi)\bigr).
\end{aligned}
\]
The until witness may be \(j=0\).  Atomic and Boolean clauses have their usual
state meanings.  Finally,
\[
 \operatorname{Sat}(\Phi)=\{s\in S:s\models\Phi\},
 \qquad
 TS\models\Phi\iff I\subseteq\operatorname{Sat}(\Phi).
\]
In particular, the system clause is vacuously true if \(I=\varnothing\); the
book explicitly permits that case \cite[defs.~6.1, 6.4--6.5;
pp.~317--321]{BaierKatoen2008}.

With these source satisfaction clauses now in place, forgetting action names
does not change any state sequence.  For nonempty \(S\), every Baier--Katoen
CTL state formula \(\Phi\), and every \(s\in S\), induction on the translation
gives
\[
 TS,s\models_{\mathrm{BK}}\Phi
 \iff
 \mathcal K_0,s\models\operatorname{tr}(\Phi).
\]
For finite nonempty \(S\), Method~\ref{def:p09-ctl-labelling} therefore
returns exactly
\[
 \{s\in S:\mathcal K_0,s\models\operatorname{tr}(\Phi)\}
 =\{s\in S:TS,s\models_{\mathrm{BK}}\Phi\},
\]
and Theorem~\ref{thm:p09-ctl-labelling} supplies correctness and the host bound
\[
 O\bigl((|S|+|R_0|)\,|\operatorname{tr}(\Phi)|\bigr).
\]

\paragraph{Existential normal form and the translation boundary.}
Definition~6.13 calls the following class existential normal form:
\[
 \Phi::=true\mid a\mid\Phi_1\land\Phi_2\mid\neg\Phi
       \mid\exists\bigcirc\Phi
       \mid\exists(\Phi_1\mathbin{U}\Phi_2)
       \mid\exists\Box\Phi.
\]
Theorem~6.14 gives rewrites eliminating universal modalities.  Its proof is a
rewrite outline, not a displayed structural induction, and the source warns
that the right operand can be duplicated enough to make an explicit
translation exponentially larger.  Consequently Algorithm~14 accepts an ENF
formula for exposition, but the complexity theorem must be implemented by
handling the original modalities directly rather than by first materializing
that translation \cite[pp.~332--333, 355]{BaierKatoen2008}.

\paragraph{Satisfaction sets and their proof.}
Algorithm~13 processes genuine subformulas in increasing length and finally
tests \(I\subseteq\operatorname{Sat}(\Phi)\).  Theorem~6.23 supplies the
set equations needed by that traversal.  Besides the atomic and Boolean
clauses, it gives
\[
 \operatorname{Sat}(\exists\bigcirc\Phi)
 =\{s:\operatorname{Post}(s)\cap\operatorname{Sat}(\Phi)\ne\varnothing\}.
\]
It characterizes \(\operatorname{Sat}(\exists(\Phi\mathbin{U}\Psi))\) as
the least \(T\subseteq S\) satisfying
\[
 \operatorname{Sat}(\Psi)\subseteq T,
 \qquad
 s\in\operatorname{Sat}(\Phi)\ \land\
 \operatorname{Post}(s)\cap T\ne\varnothing
 \ \Longrightarrow\ s\in T,
\]
and \(\operatorname{Sat}(\exists\Box\Phi)\) as the greatest
\(T\subseteq S\) satisfying
\[
 T\subseteq\operatorname{Sat}(\Phi),
 \qquad
 s\in T\ \Longrightarrow\
 \operatorname{Post}(s)\cap T\ne\varnothing.
\]
The least and greatest qualifications are with respect to inclusion.  They
are the source-level content represented by the \(\mu\)- and \(\nu\)-notation
in Method~\ref{def:p09-ctl-labelling}; that notation is not substituted
for the book's stated set characterizations.

For until, the source uses the semantic expansion law to prove the two closure
properties, then follows a finite witness backward: induction on the distance
to its final \(\Psi\)-state puts every earlier \(\Phi\)-state into any set
having those properties, proving leastness.  For \(\exists\Box\Phi\), every
satisfying path stays inside \(\operatorname{Sat}(\Phi)\), proving the two
displayed conditions.  In the converse direction the source, under its
standing no-terminal-state convention, repeatedly selects an in-set successor
to construct a path that stays in the set, proving greatestness.  These are the
detailed arguments it prints for the two nontrivial clauses; clauses (a)--(e)
of Theorem~6.23 are called straightforward rather than expanded line by line.

The handbook adds one clarification to that source proof: the displayed
post-fixed-set condition itself supplies an in-\(T\) successor at every stage.
Recursive choice therefore constructs the required infinite path; when the
state space is finite, repeated states instead expose a cycle.  This
clarification neither removes the source's standing no-terminal-state
semantic convention nor changes the stated greatest-set result
\cite[thm.~6.23; pp.~343--345]{BaierKatoen2008}.

Algorithm~14 has its own complete input/output contract.  Given a finite
transition system and an ENF formula, it returns the formula's satisfaction set
by the following seven branches:
\[
\begin{aligned}
 \operatorname{Sat}(true)&=S,\qquad&
 \operatorname{Sat}(a)&=\{s\in S:a\in L(s)\},\\
 \operatorname{Sat}(\Phi_1\land\Phi_2)
   &=\operatorname{Sat}(\Phi_1)\cap\operatorname{Sat}(\Phi_2),&
 \operatorname{Sat}(\neg\Psi)&=S\setminus\operatorname{Sat}(\Psi),\\
 \operatorname{Sat}(\exists\bigcirc\Psi)
   &=\{s:\operatorname{Post}(s)\cap\operatorname{Sat}(\Psi)
             \ne\varnothing\}.&&
\end{aligned}
\]
For \(\exists(\Phi_1\mathbin{U}\Phi_2)\), it initializes
\(T:=\operatorname{Sat}(\Phi_2)\) and repeatedly chooses one
\[
 s\in\operatorname{Sat}(\Phi_1)\setminus T
 \quad\text{with}\quad
 \operatorname{Post}(s)\cap T\ne\varnothing
\]
and adds \(s\), stopping when no such state remains.  For
\(\exists\Box\Phi\), it initializes \(T:=\operatorname{Sat}(\Phi)\) and
repeatedly chooses and deletes one \(s\in T\) for which
\(\operatorname{Post}(s)\cap T=\varnothing\).  These are the source's naive
one-state-at-a-time loops, not the optimized worklists of Algorithms~15--16.
Finiteness and Theorem~6.23 identify their results with the least and greatest
sets above \cite[alg.~14; p.~348]{BaierKatoen2008}.

Equivalently, their simultaneous approximation chains are
\[
\begin{aligned}
 T_0&=\operatorname{Sat}(\Psi),&
 T_{i+1}&=T_i\cup
   \{s\in\operatorname{Sat}(\Phi):
        \operatorname{Post}(s)\cap T_i\ne\varnothing\},\\
 U_0&=\operatorname{Sat}(\Phi),&
 U_{i+1}&=U_i\cap
   \{s\in\operatorname{Sat}(\Phi):
        \operatorname{Post}(s)\cap U_i\ne\varnothing\}.
\end{aligned}
\]
Finiteness makes both chains stabilize.  Algorithm~15 realizes the increasing
chain with a work set initially equal to \(\operatorname{Sat}(\Psi)\) and
inverse adjacency lists: each extracted state contributes previously unseen
\(\Phi\)-predecessors.  Algorithm~16 realizes the decreasing chain by putting
\(S\setminus\operatorname{Sat}(\Phi)\) in a work set, assigning each
\(s\in\operatorname{Sat}(\Phi)\) the count
\(\lvert\operatorname{Post}(s)\rvert\), and deleting a predecessor when its
count reaches zero.  Writing \(T\) for the retained set and \(E\) for the
current work set, its invariant is
\[
 \forall s\in T,\qquad
 count[s]=\lvert\operatorname{Post}(s)\cap(T\cup E)\rvert,
\]
\cite[algs.~15--16; pp.~349--352]{BaierKatoen2008}.

Now induct on the length of a subformula in Algorithm~13.  The clauses for
\(true\), atoms, negation, and conjunction compute their semantic sets from
already correct proper subformulas.  The next clause is the displayed
successor equation.  The two worklist clauses terminate at exactly the least
and greatest sets just proved.  Hence every table entry equals its semantic
satisfaction set, in particular the entry for the input formula; the final
subset test is exactly the source definition of \(TS\models\Phi\).  This
paragraph makes explicit the global induction assembled from Algorithm~13 and
Theorem~6.23; the book does not print it as a separately labelled theorem
proof.

\paragraph{The source's worked graph.}
Examples~6.26--6.28 use the eight states labelled
\[
\begin{array}{c|cccccccc}
s&s_0&s_1&s_2&s_3&s_4&s_5&s_6&s_7\\\hline
L(s)&\{a,b,c\}&\{a,b\}&\{b,c\}&\{a\}&
      \{b\}&\{a,c\}&\{c\}&\varnothing
\end{array}
\]
and the edges
\[
\begin{gathered}
s_0\to s_2,\quad s_1\to s_3,\quad
s_2\to s_0,\ s_2\to s_1,\quad s_3\to s_0,\\
s_4\to s_0,\ s_4\to s_1,\quad
s_5\to s_1,\ s_5\to s_7,\quad s_6\to s_4,\quad
s_7\to s_3,\ s_7\to s_6.
\end{gathered}
\]
Figure~6.11 marks \(I=\{s_1\}\).  That component is immaterial to the
global satisfaction-set calculations below, but it is required for the final
system-level test \(I\subseteq\operatorname{Sat}(\Phi)\).

The printed example compares the Boolean truth values of \(a,b,c\), whereas
the grammar above treats them as atomic propositions.  To make that notation
eliminable, put
\[
 \alpha\lor\beta:=\neg(\neg\alpha\land\neg\beta),\qquad
 (\alpha=\beta):=(\alpha\land\beta)\lor(\neg\alpha\land\neg\beta),
 \qquad
 (\alpha\ne\beta):=(\alpha\land\neg\beta)\lor(\neg\alpha\land\beta).
\]
These comparison signs are metalinguistic propositional abbreviations needed
only to parse the source's formula; they add neither a new state-variable sort
nor new atomic predicates.  Thus, for
\(\Theta=((a=c)\land(a\ne b))\), direct evaluation of the label table gives
\(\operatorname{Sat}(\Theta)=\{s_4,s_5\}\).  Backward search for
\(\exists\Diamond\Theta=\exists(true\mathbin{U}\Theta)\) first adds \(s_6\)
and then \(s_7\), so
\[
 \operatorname{Sat}(\exists\Diamond\Theta)
   =\{s_4,s_5,s_6,s_7\}.
\]
For \(\exists\Box b\), Algorithm~16 starts from
\(\{s_0,s_1,s_2,s_4\}\) and removes \(s_1\), returning
\(\{s_0,s_2,s_4\}\).  Example~6.28 obtains the same result by restricting
to the \(b\)-states: \(\{s_0,s_2\}\) is its only nontrivial strongly
connected component, \(s_4\) can reach it, and \(s_1\) cannot.  Thus the
examples instantiate both fixed-point directions on the source's own graph,
rather than on a replacement example \cite[pp.~349--354]{BaierKatoen2008}.

\paragraph{Exact complexity claim.}
Let \(N\) be the number of states and \(K\) the number of transitions.  With
predecessor sets stored as linked lists, Algorithms~15--16 take
\(O(N+K)\) time for one temporal subformula.  The parse-tree traversal is
linear in \(\lvert\Phi\rvert\), and a linked list of initial states makes the
last subset test linear in \(\lvert I\rvert\).  Theorem~6.30 therefore states
\[
 \operatorname{time}(TS\models\Phi)
   \in O((N+K)\lvert\Phi\rvert).
\]
This bound inherits finiteness, the no-terminal-state convention, the stated
adjacency representation, and direct treatment of modalities.  It does not
license an exponentially expanded ENF intermediate.  Although the section
heading mentions time and space, these pages print no separate asymptotic
space theorem \cite[p.~355]{BaierKatoen2008}.

The linear dependence on the supplied CTL formula is not a comparison of how
succinctly different logics express the same property.  Example~6.31 makes
that qualification explicit.  For \(G=(V,E)\) with \(n\geq1\) and
\(V=\{v_1,\ldots,v_n\}\), it first defines the following transition system
without terminal states:
\[
 TS(G)=(V\cup\{b\},\{\tau\},{\to},V,V,L),
\]
where
\[
 L(v_i)=\{v_i\},\qquad L(b)=\varnothing,
\]
and
\[
 v_i\xrightarrow{\tau}v_j\iff(v_i,v_j)\in E,
 \qquad
 x\xrightarrow{\tau}b\quad(x\in V\cup\{b\}).
\]
Thus every graph vertex is initial, every state can enter \(b\), and \(b\)
loops to itself; the added state exists solely to preserve the chapter's
no-terminal-state hypothesis.

The source then defines
\[
 \Psi(v_{i_1})=v_{i_1},
\]
and, for \(2\le k\le n\),
\[
 \Psi(v_{i_1},\ldots,v_{i_k})
   =v_{i_1}\land\exists\bigcirc
       \Psi(v_{i_2},\ldots,v_{i_k}).
\]
Let \(\Phi_n\) be the disjunction of
\(\Psi(v_{i_1},\ldots,v_{i_n})\) over all permutations of
\(1,\ldots,n\).  Direct induction on the tuple length gives
\[
 \operatorname{Sat}(\Psi(v_{i_1},\ldots,v_{i_n}))
 =
 \begin{cases}
  \{v_{i_1}\},&
    v_{i_1},\ldots,v_{i_n}\text{ is a Hamiltonian path in }G,\\
  \varnothing,&\text{otherwise}.
 \end{cases}
\]
Consequently,
\[
 G\text{ has a Hamiltonian path}
 \iff TS(G)\not\models\neg\Phi_n.
\]
The explicit enumeration makes \(\Phi_n\) exponential in \(n\), while the
book recalls a polynomial-length LTL encoding.  It informally argues that
polynomial-length equivalent CTL formulas would combine with polynomial CTL
model checking to imply \(\mathsf P=\mathsf{NP}\).  That printed inference
omits a uniformity premise: it is valid if a polynomial-time procedure
constructs the polynomial-size CTL family, whereas mere existence of such a
family gives only nonuniform advice and does not by itself imply
\(\mathsf P=\mathsf{NP}\).  This presentation records that gap rather than
promoting the source's informal argument to a theorem.  The example remains a
formula-succinctness warning, not part of Theorem~6.30 itself
\cite[ex.~6.31; pp.~355--358]{BaierKatoen2008}.

\paragraph{CTL versus LTL: the actual witness systems.}
Theorem~6.21 states incomparability, but its proof surface is not complete.
For the LTL formula \(\Diamond\Box a\), the base system \(TS_0\) has an
empty-labelled initial state \(s_0\), an \(a\)-labelled state \(t_0\), and
edges \(s_0\to t_0\), \(t_0\to t_0\), and \(t_0\to s_0\).  The base system
\(TS'_0\) instead has empty-labelled initial \(s'_0\), \(a\)-labelled
\(t'_0\), and only \(s'_0\to t'_0\), \(t'_0\to t'_0\).

For \(n>0\), both inductive members contain a copy of \(TS'_{n-1}\).
The unprimed \(TS_n\) adds fresh empty-labelled initial \(s_n\) and fresh
\(a\)-labelled \(t_n\), with exactly the new edges
\[
 s_n\to t_n,\qquad t_n\to t_n,
 \qquad t_n\to s'_{n-1},\qquad t_n\to s_n.
\]
The primed \(TS'_n\) adds fresh empty-labelled initial \(s'_n\) and fresh
\(a\)-labelled \(t'_n\), with exactly the new edges
\[
 s'_n\to t'_n,\qquad t'_n\to t'_n,
 \qquad t'_n\to s'_{n-1}.
\]
The unprimed member fails \(\Diamond\Box a\) on
\((s_nt_n)^\omega\); every path of the primed member eventually remains in
an \(a\)-labelled state. The source omits the indispensable bounded-formula
induction. The next result supplies it for this family; the source's second
LTL witness is still only said to be similar.

\begin{hbtheorem}{thm:p09-bk-ctl-witness-depth}{proved=annex:p09-bk-ctl-witness-depth}{Bounded CTL depth in the witness family}
Use the systems \(TS_n,TS'_n\) just defined, with the separate base
case at \(n=0\). For a CTL formula \(\Phi\) over \(\{a\}\), let
\(d(\Phi)\) be its temporal depth: atoms have depth zero, negation
preserves depth, a Boolean combination takes the maximum, and a paired
path-and-temporal operator adds one to the maximum depth of its state
operands. For every \(n\geq d(\Phi)\),
\[
\begin{aligned}
 TS_n,s_n\models\Phi&\quad\Longleftrightarrow\quad TS'_n,s'_n\models\Phi,\\
 TS_n,t_n\models\Phi&\quad\Longleftrightarrow\quad TS'_n,t'_n\models\Phi.
\end{aligned}
\]
Moreover, the truth values at \(s'_m\) and at \(t'_m\), respectively,
are independent of \(m\) once \(m\geq d(\Phi)\). In particular no
CTL state formula defines the universal-path LTL property
\(\Diamond\Box a\) on finite total transition structures.

This is the handbook's completion of the bounded-formula step in
Baier--Katoen's proof of Theorem~6.21(a)
\cite[pp.~337--340]{BaierKatoen2008}. Temporal depth is at most formula
length, so it also proves their stated length bound for these systems.
The other LTL witness and the separately omitted Theorem~6.18 retain
their own proof obligations.
\end{hbtheorem}


For the other direction, Figure~6.10(a) has an empty-labelled state \(s\)
with a self-loop and an edge to an \(a\)-labelled self-loop state \(s'\).  It
satisfies \(\forall\Box\exists\Diamond a\).  Figure~6.10(b) retains only the
empty-labelled self-loop and fails that CTL formula.  If an LTL formula were
equivalent to it, the first system's path \(s^\omega\) would satisfy that LTL
formula; because the same word is the only path of the second system, the
second system would satisfy it too, a contradiction.  This trace argument is
printed, but Theorem~6.18, used for two additional witnesses, has its proof
expressly omitted \cite[pp.~334--340]{BaierKatoen2008}.

\paragraph{CTL* typing, embedding, and strictness.}
Definition~6.80 replaces the CTL path grammar by
\[
\begin{aligned}
 \Phi&::=true\mid a\mid\Phi_1\land\Phi_2\mid\neg\Phi\mid\exists\varphi,\\
 \varphi&::=\Phi\mid\varphi_1\land\varphi_2\mid\neg\varphi
             \mid\bigcirc\varphi\mid\varphi_1\mathbin{U}\varphi_2.
\end{aligned}
\]
Here a state formula can be an elementary path formula, and path-level Boolean
and temporal nesting is legal.  Universal path quantification is derived by
\(\forall\varphi:=\neg\exists\neg\varphi\).  At path level,
\(\pi\models\Phi\) iff its first state satisfies \(\Phi\), next evaluates
the suffix \(\pi[1..]\), and until evaluates the appropriate suffixes at a
non-strict witness position.

For every transition system without terminal states, every state \(s\), and
every LTL formula \(\varphi\), Theorem~6.83 states
\[
 s\models_{\mathrm{LTL}}\varphi
 \quad\Longleftrightarrow\quad
 s\models_{\mathrm{CTL}^*}\forall\varphi.
\]
This is a direct match of the two semantic definitions, not a CTL* model-
checking algorithm.  Theorem~6.84 then uses the well-typed state formula
\[
 (\forall\Diamond\Box a)\lor
 (\forall\Box\exists\Diamond b)
\]
and states that it has neither an LTL nor a CTL equivalent.  Its printed proof
is one sentence referring back to Theorem~6.21.  The omitted elementary step
is substitution: make \(a\) false everywhere to isolate the CTL-only
disjunct, or make \(b\) false everywhere to isolate the LTL-only disjunct.
Thus CTL* contains both languages strictly, while the bottom-up CTL algorithm
above establishes no CTL* model-checking result
\cite[defs.~6.80--6.82 and thms.~6.83--6.84; pp.~422--425]{BaierKatoen2008}.

\paragraph{Printed defects and proof boundary.}
This presentation uses the formulas forced by the surrounding definitions
rather than copying several local defects: p.~333 drops a negation in the
expansion of \(\forall\Box\Phi\); p.~338 mislabels the primed witness tuple
and writes an inductive clause as though an undefined index \(-1\) existed; p.~340
changes \(\forall\Box\exists\Diamond a\) to a different modality in one
line; p.~347 repeats the wrong operand in the universal weak-until expansion;
p.~352 attaches a stale comment to later Algorithm~16 work items; and p.~354
misnumbers two subfigures.  Theorem~6.21 also leaves its central bounded-
formula induction to the reader, while Theorem~6.84 supplies only the cited
one-sentence argument.  The source provides a complete body and complete
algorithm descriptions, but not a complete printed proof of every theorem.
The presentation above therefore instantiates
Method~\ref{def:p09-ctl-labelling} and supports
Theorem~\ref{thm:p09-ctl-labelling} only at these explicitly separated source
and handbook-proof strengths.
